\documentclass[aps,preprint]{revtex4}%
\usepackage{amsfonts}
\usepackage{amsmath}
\usepackage{amssymb}
\usepackage{graphicx}%
\setcounter{MaxMatrixCols}{30}
%TCIDATA{OutputFilter=latex2.dll}
%TCIDATA{Version=4.10.0.2363}
%TCIDATA{CSTFile=revtex4.cst}
%TCIDATA{Created=Monday, September 20, 2004 09:58:21}
%TCIDATA{LastRevised=Tuesday, March 01, 2005 09:46:06}
%TCIDATA{<META NAME="GraphicsSave" CONTENT="32">}
%TCIDATA{<META NAME="DocumentShell" CONTENT="Articles\SW\REVTeX 4">}
%TCIDATA{Language=American English}
\newtheorem{theorem}{Theorem}
\newtheorem{acknowledgement}[theorem]{Acknowledgement}
\newtheorem{algorithm}[theorem]{Algorithm}
\newtheorem{axiom}[theorem]{Axiom}
\newtheorem{claim}[theorem]{Claim}
\newtheorem{conclusion}[theorem]{Conclusion}
\newtheorem{condition}[theorem]{Condition}
\newtheorem{conjecture}[theorem]{Conjecture}
\newtheorem{corollary}[theorem]{Corollary}
\newtheorem{criterion}[theorem]{Criterion}
\newtheorem{definition}[theorem]{Definition}
\newtheorem{example}[theorem]{Example}
\newtheorem{exercise}[theorem]{Exercise}
\newtheorem{lemma}[theorem]{Lemma}
\newtheorem{notation}[theorem]{Notation}
\newtheorem{problem}[theorem]{Problem}
\newtheorem{proposition}[theorem]{Proposition}
\newtheorem{remark}[theorem]{Remark}
\newtheorem{solution}[theorem]{Solution}
\newtheorem{summary}[theorem]{Summary}
\newenvironment{proof}[1][Proof]{\noindent\textbf{#1.} }{\ \rule{0.5em}{0.5em}}
\begin{document}
\preprint{ }
\title[Janaks Theorem]{Generalized Janaks Theorem}
\author{B. Weiner}
\affiliation{Pennsylvania State University}
\keywords{one two three}
\pacs{PACS number}

\begin{abstract}
Shell document for REV\TeX{} 4.

\end{abstract}
\date{12/21/2004}
\startpage{1}
\maketitle
\tableofcontents


\section{Introduction\label{Intro}}

\bigskip

\section{Standard Janaks Theorem\label{Standard}}

\subsection{Hohenberg-Kohn DFT}

\bigskip Density Functional Theory is based on the existence of a universal
functional, $\mathcal{\tilde{F}}\left(  \rho\right)  $ \cite{HK} of the charge
density $\rho\left(  \mathbf{r}\right)  ,$ where
\begin{equation}
\rho\left(  \mathbf{r}\right)  =\int D^{1}\left(  \mathbf{r,\xi;r,\xi}\right)
d\mathbf{\xi}%
\end{equation}
and $D^{1}\left(  \mathbf{r,\xi;r}^{\prime}\mathbf{,\xi}^{\prime}\right)  $ is
the integral kernel of the First Order Reduced Operator (FORDO) $D^{1}$
\cite{Coleman} (and see Appendix \ref{DenMat} for details of density and
reduced density operators and matrices), which can be expanded as
\begin{equation}
D^{1}=\int D^{1}\left(  \mathbf{r,\xi;r}^{\prime}\mathbf{,\xi}^{\prime
}\right)  \left\vert \mathbf{r,\xi}\right\rangle \left\langle \mathbf{r}%
^{\prime}\mathbf{,\xi}^{\prime}\right\vert d\mathbf{r}d\mathbf{r}^{\prime
}d\mathbf{\xi}d\mathbf{\xi}^{\prime}%
\end{equation}
in the eigenbasis $\left\{  \left\vert \mathbf{r,\xi}\right\rangle \right\}  $
of the position and spin operators. In order to consider spin dependent
effects it is convenient to extend the original formulation and consider a
universal functional, $\mathcal{F}\left(  \gamma\right)  $, of the space-spin
density $\gamma\left(  \mathbf{r,\xi}\right)  =$ $D^{1}\left(  \mathbf{r,\xi
;r,\xi}\right)  .$ For a system that contains on average $N$ electrons the
space spin density must satisfy the condition
\begin{equation}
\int\gamma\left(  \mathbf{r,\xi}\right)  d\mathbf{r}d\mathbf{\xi}=N
\end{equation}
and one can define equivalence classes, $\left\{  \left[  \gamma\right]
_{\left\langle N\right\rangle }\right\}  $ of state density operators by
\begin{align*}
\left[  \gamma\right]  _{\left\langle N\right\rangle }  &  =\left\{  \left.
D\right\vert \Xi^{0}\left(  D\right)  =\gamma;D\in S_{\left\langle
N\right\rangle }\right\} \\
S_{\left\langle N\right\rangle }  &  =\left\{  \left.  D\right\vert
D\geq0;\operatorname{Tr}D=N\right\} \\
D  &  :\text{Fermi Fock Space }\rightarrow\text{Fermi Fock Space}%
\end{align*}%
\[
\Phi\left(  \mathbf{x}\right)  =\sum_{1\leq j\leq2s}\varphi_{j}\left(
\mathbf{x}\right)  ^{\ast}a_{j}%
\]%
\[
\gamma\left(  \mathbf{x}\right)  =\operatorname{Tr}\left\{  \Phi\left(
\mathbf{x}\right)  ^{\dagger}\Phi\left(  \mathbf{x}\right)  D\right\}
\]
(see Appendix \ref{DenMat} for the definitions of the map $\Xi^{0}$) and hence
of $2^{nd}$ order reduced density operators by
\[
\left[  \gamma\right]  _{\left\langle N\right\rangle 2}=\left\{  \left.
D^{2}\right\vert D^{2}=\binom{N}{2}L_{2}\left(  D\right)  \in\left[
\gamma\right]  _{\left\langle N\right\rangle }\right\}  .
\]


Even though the explicit form of the functional $\mathcal{F}_{HK}\left(
\gamma\right)  $ is not known one can define it implicitly \cite{LevyLieb}
\begin{align*}
\mathcal{F}_{HK}\left(  \gamma\right)   &  =\operatorname{Tr}\left\{
\overset{}{KD_{\ast}^{1}\left(  \gamma\right)  +V_{2}D_{\ast}^{2}\left(
\gamma\right)  }\right\} \\
&  =\min_{\substack{D^{2}\in\left[  \mathbf{\gamma}\right]  _{\left\langle
N\right\rangle 2}\\D^{1}=\frac{2}{N-1}L_{2}^{1}\left(  D^{2}\right)
}}\operatorname{Tr}\left\{  \overset{}{KD^{1}+V_{2}D^{2}}\right\}  ,
\end{align*}
where the kinetic energy operator $K$ is
\begin{equation}
K=\int\left\langle \mathbf{r}^{\prime}\mathbf{,\xi}\left\vert \left(
-\frac{1}{2}\nabla_{\mathbf{r}}^{2}\right)  \mathbf{r,\xi}\right.
\right\rangle \left\vert \mathbf{r,\xi}\right\rangle \left\langle
\mathbf{r}^{\prime}\mathbf{,\xi}\right\vert d\mathbf{r}d\mathbf{r}^{\prime
}d\mathbf{\xi}%
\end{equation}
and the coulomb operator
\begin{equation}
V_{2}=\frac{1}{4}\int\frac{1}{\left\Vert \mathbf{r-r}^{\prime}\right\Vert
}\left\vert \mathbf{r,\mathbf{r}^{\prime},\xi,\xi}^{\prime}\right\rangle
\left\langle \mathbf{r,\mathbf{r}^{\prime},\xi,\xi}^{\prime}\right\vert
d\mathbf{r}d\mathbf{r}^{\prime}d\mathbf{\xi}d\mathbf{\xi}^{\prime}.
\end{equation}
This allows one to define kinetic energy, coulomb and exchange - correlation
functionals on paths of representable \cite{Kummer2} Second Order Reduced
Density Operators (SORDO's)$\ D_{\ast}^{2}\left(  \gamma\right)  $ given by
eq. $\left(  \ref{HK}\right)  $ as
\begin{align*}
\mathcal{F}_{K}\left(  \gamma\right)   &  =\operatorname{Tr}\left\{  KD_{\ast
}^{1}\left(  \gamma\right)  \right\}  ,\\
\mathcal{F}_{C}\left(  \gamma\right)   &  =\int\frac{\gamma\left(
\mathbf{r,\xi}\right)  \gamma\left(  \mathbf{r}^{\prime}\mathbf{,\xi}\right)
}{\left\Vert \mathbf{r-r}^{\prime}\right\Vert }d\mathbf{r}d\mathbf{r}^{\prime
}d\mathbf{\xi},\\
\mathcal{F}_{XC}\left(  \gamma\right)   &  =\operatorname{Tr}\left\{
V_{2}D_{\ast}^{2}\left(  \gamma\right)  \right\}  -\mathcal{F}_{C}\left(
\gamma\right)  ,
\end{align*}
producing the decomposition
\[
\mathcal{F}_{HK}\left(  \gamma\right)  =\mathcal{F}_{K}\left(  \gamma\right)
+\mathcal{F}_{C}\left(  \gamma\right)  +\mathcal{F}_{XC}\left(  \gamma\right)
.
\]
$\lceil$%
\begin{align}
\mathcal{F}_{HK}\left(  \alpha_{1}\gamma_{1}+\alpha_{2}\gamma_{2}\right)   &
=\alpha_{1}\min_{\substack{D^{2}\in\left[  \mathbf{\gamma}_{1}\right]
_{\left\langle N\right\rangle 2}\\D^{1}=\frac{2}{N-1}L_{2}^{1}\left(
D^{2}\right)  }}\operatorname{Tr}\left\{  \overset{}{\left\{  KD^{1}%
+V_{2}D^{2}\right\}  }\right\}  +\alpha_{2}\min_{\substack{D^{2}\in\left[
\mathbf{\gamma}_{2}\right]  _{\left\langle N\right\rangle 2}\\D^{1}=\frac
{2}{N-1}L_{2}^{1}\left(  D^{2}\right)  }}\operatorname{Tr}\left\{  \overset
{}{\left\{  KD^{1}+V_{2}D^{2}\right\}  }\right\} \nonumber\\
&  =\alpha_{1}\mathcal{F}_{HK}\left(  \gamma_{1}\right)  +\alpha
_{2}\mathcal{F}_{HK}\left(  \gamma_{2}\right)
\end{align}
Thus $\mathcal{F}_{HK}$ is a convex function.$\rfloor$ The total energy
functional of a system that contains on average $N$ electrons and experiences
an environment described by a \emph{local} potential $\mathrm{v}$ is
\[
\mathcal{E}\left(  \gamma\mathbf{,}\mathrm{v}\right)  =\mathcal{F}_{HK}\left(
\gamma\right)  +f_{\mathrm{v}}\left(  \gamma\right)  ,
\]
where the explicit effect of the environment is described by the functional
\[
f_{\mathrm{v}}\left(  \gamma\right)  =\int\mathrm{v}\left(  \mathbf{r,\xi
}\right)  \gamma\left(  \mathbf{r,\xi}\right)  d\mathbf{r}d\mathbf{\xi.}%
\]
The ground state energy $E_{GS}\left(  \mathrm{v}\right)  $ and space-spin
density $\gamma_{GS}$ $\left(  \mathrm{v}\right)  $ are then determined by
\begin{equation}
E_{GS}\left(  \mathrm{v}\right)  =\mathcal{E}\left(  \gamma_{GS}%
\mathbf{,}\mathrm{v}\right)  =\min_{\mathbf{\gamma\in}\mathcal{D}}%
\mathcal{E}\left(  \gamma\mathbf{,}\mathrm{v}\right)  ,
\end{equation}
where
\begin{equation}
\mathcal{D=}\left\{  \gamma\left\vert \int\gamma\left(  \mathbf{r,\xi}\right)
d\mathbf{r}d\mathbf{\xi}=N,\quad\gamma\left(  \mathbf{r,\xi}\right)
\geq0\right.  \right\}  .
\end{equation}
$\lceil$The domain $\mathcal{D}$ is a compact (closed in $L_{1}$ topology and
bounded) convex set thus a unique minimum exists$\rfloor$

The density $\gamma_{GS}$ $\left(  \mathrm{v}\right)  $ can be always
expressed in terms of the Natural General Spin Orbitals (NGSO's) $\left\{
\varphi_{j}\right\}  $ and occupation numbers $\left\{  N_{j}\right\}  $ of
$D_{\ast}^{1}\left(  \gamma\right)  $ as%

\begin{equation}
\gamma_{GS}\left(  \mathrm{v}\right)  \left(  \mathbf{x}\right)  =D_{\ast}%
^{1}\left(  \gamma\right)  \left(  \mathbf{x,x}\right)  =\sum_{1\leq j\leq
r}N_{j}\left\vert \varphi_{j}\left(  \mathbf{x}\right)  \right\vert ^{2},
\end{equation}
but such expansions are not unique, in particular Harimanns \cite{Harimann}
theorem ensures that there are always expansions of the form
\begin{equation}
\gamma_{GS}\left(  \mathrm{v}\right)  \left(  \mathbf{x}\right)  =D_{\ast}%
^{1}\left(  \gamma\right)  \left(  \mathbf{x,x}\right)  =\sum_{1\leq j\leq
N}\left\vert \varphi_{j}\left(  \mathbf{x}\right)  \right\vert ^{2}%
\end{equation}
in terms of square summable functions $\left\{  \left.  \varphi_{j}\right\vert
1\leq j\leq N\right\}  $.

\subsection{Kohn-Sham DFT}

\bigskip One can introduce a universal kinetic energy functional
\[
\mathcal{F}_{T}\left(  \gamma\right)  =\operatorname{Tr}\left\{  KD_{\#}%
^{1}\right\}  =\min_{D^{1}\in\left[  \gamma\right]  _{\left\langle
N\right\rangle 1}}\operatorname{Tr}\left\{  KD^{1}\right\}  ,
\]


$\lceil$As $K$ is not a bounded operator (thus does not define a continuous
functional of $D^{1}$) the $\min$ might not exist and should be replaced with
$\inf$ even though the domain $\left[  \gamma\right]  _{\left\langle
N\right\rangle 1}$ is a compact convex set. The resulting functional is
convex
\begin{equation}
\mathcal{F}_{T}\left(  \gamma\right)  =\mathcal{F}_{T}\left(  \alpha_{1}%
\gamma_{1}+\alpha_{2}\gamma_{2}\right)  =\alpha_{1}\mathcal{F}_{T}\left(
\gamma_{1}\right)  +\alpha_{2}\mathcal{F}_{T}\left(  \gamma_{2}\right)
\end{equation}
$\rfloor$ which can be expanded as
\begin{equation}
\mathcal{F}_{T}\left(  \gamma\right)  =\sum_{1\leq j\leq r}\left\langle
\varphi_{\#j}\left\vert K\varphi_{\#j}\right.  \right\rangle N_{\#j},
\end{equation}
where
\begin{align}
D_{\#}^{1}\left(  \mathbf{x,x}^{\prime}\right)   &  =\sum_{1\leq j\leq r}%
N_{j}\varphi_{\#j}\left(  \mathbf{x}\right)  \varphi_{\#j}\left(
\mathbf{x}^{\prime}\right)  ^{\ast}\nonumber\\
\gamma\left(  \mathbf{x}\right)   &  =\sum_{1\leq j\leq r}\left\vert
\varphi_{\#j}\left(  \mathbf{x}\right)  \right\vert ^{2}N_{\#j}.
\label{gammaKS}%
\end{align}
For a given space-spin density $\gamma$ one can define a total energy
functional
\begin{equation}
\mathcal{E}_{NI}\left(  \gamma\mathbf{,}\mathrm{w}\right)  =\mathcal{F}%
_{T}\left(  \gamma\right)  +f_{\gamma}\left(  \mathrm{w}\right)
\end{equation}
that describes groups of on average $N$ noninteracting electrons in an
environment characterized by a local potential $\mathrm{w.}$ One can then
determine a local potential $\mathrm{w}_{\#}$ such that
\begin{equation}
E_{NIGS}\left(  \mathrm{w}_{\#}\right)  =\mathcal{E}_{NIGS}\left(  \gamma
_{GS}\left(  \mathrm{v}\right)  \mathbf{,}\mathrm{w}_{\#}\right)
=\min_{\mathbf{\gamma\in}\mathcal{D}}\mathcal{E}_{NI}\left(  \gamma
\mathbf{,}\mathrm{w}_{\#}\right)  .
\end{equation}
As the Hamiltonian describing noninteracting electrons is a one body
Hamiltonian the ground state density operator $D_{\#}\left(  \gamma
_{GS}\left(  \mathrm{v}\right)  \mathbf{,}\mathrm{w}_{\#}\right)  $ must be an
ensemble or a pure Independent Particle State (IPS)\ density operator. Thus
the space-spin density of the ground state of the interacting system in an
local external potential $\mathrm{v}$ is the same as space-spin density of the
ground state of a non interacting system determined by an environment
$\mathrm{w}_{\#}\mathrm{.}$

One can define a Kohn-Sham energy functional, $\mathcal{F}_{KS},$ for the
interacting system as
\[
\mathcal{F}_{KS}\left(  \gamma\right)  =\mathcal{F}_{T}\left(  \gamma\right)
+\mathcal{F}_{C}\left(  \gamma\right)  +\mathcal{F}_{KSXC}\left(
\gamma\right)  ,
\]
where the Kohn-Sham exchange functional is defined as
\[
\mathcal{F}_{KSXC}\left(  \gamma\right)  =\mathcal{F}_{XC}\left(
\gamma\right)  +\mathcal{F}_{K}\left(  \gamma\right)  -\mathcal{F}_{T}\left(
\gamma\right)  .
\]
The motivation for the introduction of $\mathcal{F}_{KS}$ is that one can
parameterize $\gamma$ in terms of $\left\{  \varphi_{\#j},N_{\#j}\right\}  $
(which are unique up to unitary transformation that mix orbitals that
correspond to degenerate occupation numbers ) through eq. $\left(
\ref{KEFunct}\right)  $ i.e.
\[
\gamma\left(  \mathbf{x}\right)  =\sum_{1\leq j\leq r}\left\vert \varphi
_{\#j}\left(  \mathbf{x}\right)  \right\vert ^{2}N_{\#j}%
\]
where
\begin{subequations}
\label{Con}%
\begin{align*}
&  \sum_{1\leq j\leq r}N_{\#j}=N\\
&  0\leq N_{\#j}\leq1;\quad1\leq j\leq r\\
&  \int\varphi_{\#j}\left(  \mathbf{x}\right)  \varphi_{\#k}\left(
\mathbf{x}\right)  ^{\ast}d\mathbf{x}=\delta_{jk};\quad1\leq j,k\leq
r\mathbf{.}%
\end{align*}
The Kohn-Sham energy functional can then be expressed in terms of the
variables $\left\{  \left.  \varphi_{\#j},N_{\#j}\right\vert 1\leq j\leq
r\right\}  $ i.e.
\end{subequations}
\[
\mathcal{F}_{KS}\left(  \mathbf{\varphi}_{\#}\mathbf{,N}_{\#}\right)
=\mathcal{F}_{KS}\left(  \gamma\left(  \mathbf{\varphi}_{\#}\mathbf{,N}%
_{\#}\right)  \right)
\]%
\begin{align*}
&  \qquad\qquad\qquad\qquad\mathcal{F}_{KS}\left(  \mathbf{\varphi}%
_{\#}\mathbf{,N}_{\#}\right) \\
&  \qquad\qquad\qquad\qquad\qquad\left.  \shortparallel\right. \\
&  \mathcal{F}_{T}\left(  \mathbf{\varphi}_{\#}\mathbf{,N}_{\#}\right)
+\mathcal{F}_{C}\left(  \mathbf{\varphi}_{\#}\mathbf{,N}_{\#}\right)
+\mathcal{F}_{KSXC}\left(  \mathbf{\varphi}_{\#}\mathbf{,N}_{\#}\right)
\end{align*}
and the minimization problem that determines the ground state energy and
space-spin density as
\[
E_{GS}=\min_{\left(  \mathbf{\varphi}_{\#}\mathbf{,N}_{\#}\right)
\mathbf{\in}\Lambda}\mathcal{E}_{KS}\left(  \mathbf{\varphi}_{\#}%
\mathbf{,N}_{\#}\mathbf{,}\mathrm{v}\right)
\]%
\[
\Lambda=\left\{  \mathbf{\varphi}_{\#}\mathbf{,N}_{\#}\right\}
\]%
\[
\mathcal{F}_{T}\left(  \mathbf{\varphi}_{\#}\mathbf{,N}_{\#}\right)
=\min_{D^{1}\left(  \mathbf{\varphi,N}\right)  \in\left[  \gamma\right]
_{\left\langle N\right\rangle 1}}\operatorname{Tr}\left\{  \sum_{1\leq j\leq
r}\left\langle \varphi_{j}\left\vert K\varphi_{j}\right.  \right\rangle
N_{j}\right\}
\]%
\begin{equation}
\mathcal{F}_{T}\left(  \gamma\right)  =\sum_{1\leq j\leq r}\left\langle
\varphi_{\#j}\left\vert K\varphi_{\#j}\right.  \right\rangle N_{\#j},
\end{equation}
where
\[
\mathcal{E}_{KS}\left(  \mathbf{\varphi}_{\#}\mathbf{,N}_{\#}\mathbf{,}%
\mathrm{v}\right)  =\mathcal{F}_{KS}\left(  \mathbf{\varphi}_{\#}%
\mathbf{,N}_{\#}\right)  +f_{\mathrm{v}}\left(  \mathbf{\varphi}%
_{\#}\mathbf{,N}_{\#}\right)
\]
and $\mathcal{\Lambda}$ is the set of feasible values for $\left(
\mathbf{\varphi}_{\#}\mathbf{,N}_{\#}\right)  $ given by eqs. $\left(
\ref{Con}\right)  .$ The Kohn-Sham constrained minimization problem can be
expressed in terms of the Lagrangian
\begin{align*}
\qquad &  \qquad\qquad\qquad\mathcal{L}\left(  \mathbf{\varphi}_{\#}%
\mathbf{,N}_{\#},\mu,\mathbf{\varepsilon,\eta,}\mathrm{v}\right) \\
&  \qquad\qquad\qquad\qquad\qquad\left.  \shortparallel\right. \\
&  \mathcal{E}_{KS}\left(  \mathbf{\varphi}_{\#}\mathbf{,N}_{\#}%
\mathbf{,}\mathrm{v}\right)  +\sum_{1\leq j\leq r}\eta_{j}\left(
1-N_{\#j}\right)  -\sum_{1\leq j\leq r}\xi_{j}N_{\#j}\\
&  -\sum_{1\leq j,k\leq r}\varepsilon_{jk}\left(  \left\langle \left.
\varphi_{\#j}\right\vert \varphi_{\#k}\right\rangle -\delta_{jk}\right)
+\mu\left(  \sum_{1\leq j\leq r}N_{\#j}-N\right)
\end{align*}
This leads to the necessary stationarity conditions
\begin{subequations}
\begin{align*}
\frac{\partial\mathcal{L}\left(  \mathbf{\varphi}_{\#}\mathbf{,N}_{\#}%
,\mu,\mathbf{\varepsilon,\eta,}\mathrm{v}\right)  }{\partial\varphi
_{\#j}^{\ast}}  &  =0\\
\frac{\partial\mathcal{L}\left(  \mathbf{\varphi}_{\#}\mathbf{,N}_{\#}%
,\mu,\mathbf{\varepsilon,\eta,}\mathrm{v}\right)  }{\partial\varphi_{\#j}}  &
=0\\
\frac{\partial\mathcal{L}\left(  \mathbf{\varphi}_{\#}\mathbf{,N}_{\#}%
,\mu,\mathbf{\varepsilon,\eta,}\mathrm{v}\right)  }{\partial N_{\#j}}  &  =0.
\end{align*}
in particular implies that
\end{subequations}
\[
\frac{\partial\mathcal{E}_{KS}\left(  \mathbf{\varphi}_{\#}\mathbf{,N}%
_{\#}\mathbf{,}\mathrm{v}\right)  }{\partial\varphi_{\#j}^{\ast}}-\sum_{1\leq
k\leq r}\varepsilon_{jk}\left\vert \varphi_{\#k}\right\rangle =0
\]
One can develope the Kohn-Sham energy functional as%
\begin{align}
&  \mathcal{E}_{KS}\left(  \mathbf{\varphi}_{\#}\mathbf{,N}_{\#}%
\mathbf{,}\mathrm{v}\right)  =\mathcal{F}_{KS}\left(  \mathbf{\varphi}%
_{\#}\mathbf{,N}_{\#}\right)  +f_{\mathrm{v}}\left(  \mathbf{\varphi}%
_{\#}\mathbf{,N}_{\#}\right) \nonumber\\
&  =\mathcal{F}_{T}\left(  \mathbf{\varphi}_{\#}\mathbf{,N}_{\#}\right)
+\mathcal{F}_{C}\left(  \mathbf{\varphi}_{\#}\mathbf{,N}_{\#}\right)
+\mathcal{F}_{KSXC}\left(  \mathbf{\varphi}_{\#}\mathbf{,N}_{\#}\right)
+f_{\mathrm{v}}\left(  \mathbf{\varphi}_{\#}\mathbf{,N}_{\#}\right)
\nonumber\\
&  =\sum_{1\leq j\leq r}\left\langle \varphi_{\#j}\left\vert K\varphi
_{\#j}\right.  \right\rangle N_{\#j}+\int\frac{\gamma\left(  \mathbf{r,\xi
}\right)  \gamma\left(  \mathbf{r}^{\prime}\mathbf{,\xi}\right)  }{\left\Vert
\mathbf{r-r}^{\prime}\right\Vert }d\mathbf{r}d\mathbf{r}^{\prime}d\mathbf{\xi
}+\int\mathrm{v}\left(  \mathbf{r,\xi}\right)  \gamma\left(  \mathbf{r,\xi
}\right)  d\mathbf{r}d\mathbf{\xi}+\mathcal{F}_{KSXC}\left(  \mathbf{\varphi
}_{\#}\mathbf{,N}_{\#}\right)  ,
\end{align}
where the minimization over the kinetic functional gives the correspondence
\begin{equation}
\gamma\left(  \mathbf{x}\right)  =\sum_{1\leq j\leq r}\left\vert \varphi
_{\#j}\left(  \mathbf{x}\right)  \right\vert ^{2}N_{\#j}.
\end{equation}
Using
\[
\frac{\partial}{\partial N_{\#j}}=\int\frac{\partial\gamma\left(
\mathbf{x}\right)  }{\partial N_{\#j}}\frac{\delta}{\delta\gamma\left(
\mathbf{x}\right)  }d\mathbf{x=}\int\varphi_{\#j}\left(  \mathbf{x}\right)
\varphi_{\#j}\left(  \mathbf{x}\right)  ^{\ast}\frac{\delta}{\delta
\gamma\left(  \mathbf{x}\right)  }d\mathbf{x}%
\]
one obtains Janaks Theorem
\begin{align*}
&  \frac{\partial\mathcal{E}_{KS}\left(  \mathbf{\varphi}_{\#}\mathbf{,N}%
_{\#}\mathbf{,}\mathrm{v}\right)  }{\partial N_{\#j}}=\int\frac{\delta
\mathcal{E}_{KS}\left(  \mathbf{\varphi}_{\#}\mathbf{,N}_{\#}\mathbf{,}%
\mathrm{v}\right)  }{\delta\gamma\left(  \mathbf{x}\right)  }\frac
{\partial\gamma\left(  \mathbf{x}\right)  }{\partial N_{\#j}}d\mathbf{x}\\
&  =\int\frac{\delta\mathcal{E}_{KS}\left(  \mathbf{\varphi}_{\#}%
\mathbf{,N}_{\#}\mathbf{,}\mathrm{v}\right)  }{\delta\gamma\left(
\mathbf{x}\right)  }\varphi_{\#j}\left(  \mathbf{x}^{\prime}\right)
\varphi_{\#j}\left(  \mathbf{x}^{\prime}\right)  ^{\ast}\frac{\delta
\gamma\left(  \mathbf{x}\right)  }{\delta\gamma\left(  \mathbf{x}^{\prime
}\right)  }d\mathbf{x}d\mathbf{x}^{\prime}\\
&  =\int\frac{\delta\mathcal{E}_{KS}\left(  \mathbf{\varphi}_{\#}%
\mathbf{,N}_{\#}\mathbf{,}\mathrm{v}\right)  }{\delta\gamma\left(
\mathbf{x}\right)  }\varphi_{\#j}\left(  \mathbf{x}^{\prime}\right)
\varphi_{\#j}\left(  \mathbf{x}^{\prime}\right)  ^{\ast}\delta\left(
\mathbf{x-x}^{\prime}\right)  d\mathbf{x}d\mathbf{x}^{\prime}\\
&  =\int\frac{\delta\mathcal{E}_{KS}\left(  \mathbf{\varphi}_{\#}%
\mathbf{,N}_{\#}\mathbf{,}\mathrm{v}\right)  }{\delta\gamma\left(
\mathbf{x}^{\prime}\right)  }\varphi_{\#j}\left(  \mathbf{x}^{\prime}\right)
\varphi_{\#j}\left(  \mathbf{x}^{\prime}\right)  ^{\ast}d\mathbf{x}^{\prime}\\
&  =\left\langle \varphi_{\#j}\left\vert \frac{\delta\mathcal{E}_{KS}\left(
\mathbf{\varphi}_{\#}\mathbf{,N}_{\#}\mathbf{,}\mathrm{v}\right)  }%
{\delta\gamma\left(  \mathbf{x}^{\prime}\right)  }\right.  \varphi
_{\#j}\right\rangle =\varepsilon_{jj}%
\end{align*}%
\begin{align*}
&  \frac{\partial\mathcal{E}_{KS}\left(  \mathbf{\varphi}_{\#}\mathbf{,N}%
_{\#}\mathbf{,}\mathrm{v}\right)  }{\partial N_{\#j}}=\int\frac{\delta
\mathcal{E}_{KS}\left(  \mathbf{\varphi}_{\#}\mathbf{,N}_{\#}\mathbf{,}%
\mathrm{v}\right)  }{\delta\gamma\left(  \mathbf{x}\right)  }\frac
{\partial\gamma\left(  \mathbf{x}\right)  }{\partial N_{\#j}}d\mathbf{x}\\
&  =\int\frac{\delta\mathcal{E}_{KS}\left(  \mathbf{\varphi}_{\#}%
\mathbf{,N}_{\#}\mathbf{,}\mathrm{v}\right)  }{\delta\gamma\left(
\mathbf{x}^{\prime}\right)  }\varphi_{\#j}\left(  \mathbf{x}^{\prime}\right)
\varphi_{\#j}\left(  \mathbf{x}^{\prime}\right)  ^{\ast}d\mathbf{x}^{\prime}\\
&  =\left\langle \varphi_{\#j}\left\vert \frac{\delta\mathcal{E}_{KS}\left(
\mathbf{\varphi}_{\#}\mathbf{,N}_{\#}\mathbf{,}\mathrm{v}\right)  }%
{\delta\gamma\left(  \mathbf{x}^{\prime}\right)  }\right.  \varphi
_{\#j}\right\rangle \\
&  \qquad\qquad\qquad\quad=\varepsilon_{jj}%
\end{align*}%
\[
\frac{\partial\mathcal{E}_{KS}\left(  \mathbf{\varphi}_{\#}\mathbf{,N}%
_{\#}\mathbf{,}\mathrm{v}\right)  }{\partial N_{\#j}}=\varepsilon_{jj}%
\]


\section{AGP Energy Functionals\label{AGP}}

\subsection{Canonical expansion of an AGP state}

\bigskip The $2N$-electron AGP state, $\left\vert g^{N}\right\rangle $, is the
$N^{th}$ antisymmetric power of the $2$-electron state described by the
geminal
\begin{equation}
\left\vert g\right\rangle =\sum_{1\leq j\leq s}n_{j}^{\frac{1}{2}}\left\vert
\varphi_{j}\varphi_{j+s}\right\rangle , \label{geminal}%
\end{equation}
where $\left\{  \left.  n_{j}^{\frac{1}{2}}\right\vert 1\leq j\leq s\right\}
$ are the positive square roots of $\left\{  \left.  n_{j}\right\vert 1\leq
j\leq s\right\}  ,$ and is given by \cite{Weiner&ortiz04}
\begin{equation}
\left\vert g^{\frac{N}{2}}\right\rangle =\left(  S_{\frac{N}{2}}\right)
^{-\frac{1}{2}}\sum_{1\leq j_{1}<\cdots<j_{\frac{N}{2}}\leq s}n_{j_{1}}%
^{\frac{1}{2}}\cdots n_{j_{s}}^{\frac{1}{2}}\left\vert \varphi_{j_{1}}%
\varphi_{j_{1}+s}\cdots\varphi_{j_{\frac{N}{2}}}\varphi_{j_{\frac{N}{2}}%
+s}\right\rangle . \label{agp}%
\end{equation}
The normalization factor, $\left(  S_{\frac{N}{2}}\right)  ^{-\frac{1}{2}}$,
is given in terms of the elementary symmetric function
\begin{equation}
S_{\frac{N}{2}}=\sum_{1\leq j_{1}<\cdots<j_{\frac{N}{2}}\leq s}n_{j_{1}}\cdots
n_{j_{N}},
\end{equation}
and the geminals $\left\vert g\right\rangle $ represent normalized
two--electron states so that
\begin{equation}
\sum_{1\leq j\leq s}n_{j}=1.
\end{equation}
It is evident from eq.\ $\left(  \ref{geminal}\right)  $ that the geminal and
geminal power state are at the least invariant to the group
\begin{equation}
SU\left(  2\right)  ^{s}=\overset{s\text{-factors}}{\overbrace{\,SU\left(
2\right)  \times\cdots\times SU\left(  2\right)  }},
\end{equation}
that is formed by $2\times2$ \emph{special} unitary transformations of the
paired CGSOs $\left\{  \left.  \left\vert \varphi_{j}\right\rangle ,\left\vert
\varphi_{j+s}\right\rangle \right\vert 1\leq j\leq s\right\}  ,$ while if
degeneracies exist amongst the $\left\{  \left.  n_{j}\right\vert 1\leq j\leq
s\right\}  $ this invariance group is larger \cite{AGPCoherent}. The canonical
CSGOs are unique up to transformations belonging this invariance group.

\subsection{Group Parameteriztion of Geminals}

Let
\[
\left\vert g_{0}\right\rangle =\sum_{1\leq j\leq s}c_{0j}\left\vert
\varphi_{0j}\varphi_{0j+s}\right\rangle ;\quad\left\Vert \left\vert
g_{0}\right\rangle \right\Vert =\sum_{1\leq j\leq s}n_{0j}=1
\]
be a reference geminal such that $c_{j}\neq0$ $\forall j$ then any geminal
$\left\vert g\left(  \mathbf{\alpha},\mathbf{\lambda},\mathbf{\theta}\right)
\right\rangle $ can be generated by the transformation
\[
\left\vert g\left(  \mathbf{\alpha},\mathbf{\lambda},\mathbf{\theta}\right)
\right\rangle =u\left(  \mathbf{\alpha}\right)  v\left(  \mathbf{\eta}\right)
w\left(  \mathbf{\theta}\right)  \left\vert g_{0}\right\rangle
\]
where
\[
u\left(  \mathbf{\alpha}\right)  =\exp i\left\{  \sum_{\substack{1\leq
j,k\leq2s\\j\neq k,k\pm s}}\alpha_{jk}a_{j}^{\dagger}a_{k}\right\}
;\quad\alpha_{jk}=\alpha_{kj}^{\ast}%
\]
%

\[
v\left(  \mathbf{\eta}\right)  =\exp\frac{1}{2}\left\{  \sum_{1\leq j\leq
s}\eta_{j}\left(  a_{j}^{\dagger}a_{j}+a_{j+s}^{\dagger}a_{j+s}\right)
\right\}
\]%
\[
w\left(  \mathbf{\theta}\right)  =\exp\frac{i}{2}\left\{  \sum_{1\leq j\leq
s}\theta_{j}\left(  a_{j}^{\dagger}a_{j}+a_{j+s}^{\dagger}a_{j+s}\right)
\right\}
\]%
\[
\left[  D_{0}^{1}\right]  _{2}=\left\{  \left.  w\left(  \mathbf{\theta
}\right)  \left\vert g_{0}\right\rangle \right\vert 0\leq\theta<2\pi\right\}
\]


Same NGSO's $\left\{  \left.  \left\vert \varphi_{j}\right\rangle \right\vert
1\leq j\leq2s\right\}  $%
\[
\left.  {}\right.
\]


different CGSO's $\left\{  \left.  \left\vert \varphi_{j}\left(
\mathbf{\theta}\right)  \right\rangle \right\vert 1\leq j\leq2s\right\}  $

New CGSO's are generated by
\begin{align*}
&  \left\vert \varphi_{j}\left(  \mathbf{\lambda}\right)  \right\rangle
=u\left(  \mathbf{\alpha}\right)  w\left(  \mathbf{\theta}\right)  \left\vert
\varphi_{0j}\right\rangle \quad\\
&  \mathbf{\lambda\equiv}\left\{  \mathbf{\alpha,\theta}\right\}
\mathbf{,\,}1\leq j\leq2s
\end{align*}
and new canonical coefficients corresponding to normalized geminals by%

\begin{align*}
c_{j}\left(  \mathbf{\eta}\right)   &  =c_{0j}e^{\eta_{j}}\mathit{s}\left(
\mathbf{\eta}\right)  ^{-\frac{1}{2}}\\
&  =p\left(  \mathbf{\eta}\right)  ^{\frac{1}{2}};\mathbf{\,}1\leq j\leq s
\end{align*}
leading to
\[
n_{j}\left(  \mathbf{\eta}\right)  =e^{2\eta_{j}}n_{0j}\mathit{s}\left(
\mathbf{\eta}\right)  ^{-1};\mathbf{\,}1\leq j\leq s.
\]


We can express the geminal in terms of unnormalized but orthogonal CGSO's as%

\begin{align*}
\left\vert g\left(  \mathbf{\Lambda}\right)  \right\rangle  &  =\sum_{1\leq
j\leq s}\left\vert \Upsilon_{j}\left(  \mathbf{\Lambda}\right)  \Upsilon
_{j+s}\left(  \mathbf{\Lambda}\right)  \right\rangle \\
&  =T\left(  \mathbf{\Lambda}\right)  \sum_{1\leq j\leq s}\left\vert
\Upsilon_{0j}\Upsilon_{0j+s}\right\rangle \\
&  =u\left(  \mathbf{\alpha}\right)  v\left(  \mathbf{\eta}\right)  w\left(
\mathbf{\theta}\right)  \sum_{1\leq j\leq s}\left\vert \Upsilon_{0j}%
\Upsilon_{0j+s}\right\rangle
\end{align*}
where
\[
\left\vert \Upsilon_{0j}\right\rangle =c_{0j}^{\frac{1}{2}}\left\vert
\varphi_{0j}\right\rangle ;\,1\leq j\leq2s
\]
and
\[
\mathbf{\Lambda\equiv}\left\{  \mathbf{\alpha,\eta,\theta}\right\}
\]
so that
\begin{align*}
\left\langle \left.  \Upsilon_{j}\left(  \mathbf{\Lambda}\right)  \right\vert
\Upsilon_{k}\left(  \mathbf{\Lambda}\right)  \right\rangle  &  =c_{j}\left(
\mathbf{\eta}\right)  \delta_{jk};\,1\leq j\leq2s\\
c_{j+s}\left(  \mathbf{\eta}\right)   &  =c_{j}\left(  \mathbf{\eta}\right)
;\,1\leq j\leq s
\end{align*}%
\[
\left\langle \left.  \Upsilon_{j}\left(  \mathbf{\Lambda}\right)  \right\vert
\Upsilon_{k}\left(  \mathbf{\Lambda}\right)  \right\rangle =c_{j}\left(
\mathbf{\eta}\right)  \delta_{jk}=e^{\eta_{j}}\mathit{s}\left(  \mathbf{\eta
}\right)  ^{-\frac{1}{2}}c_{0j}\delta_{jk}%
\]%
\[
\mathbf{n}_{j}^{\frac{1}{2}}\left(  \mathbf{\eta}\right)  =c_{0j}e^{\eta_{j}%
}\mathit{s}\left(  \mathbf{\eta}\right)  ^{-\frac{1}{2}}=c_{j}\left(
\mathbf{\eta}\right)
\]%
\begin{align*}
\left\vert \Upsilon_{j}\left(  \mathbf{\Lambda}\right)  \right\rangle  &
=c_{j}^{\frac{1}{2}}\left(  \mathbf{\eta}\right)  \left\vert \varphi
_{j}\left(  \mathbf{\lambda}\right)  \right\rangle \\
\left\Vert \Upsilon_{j}\left(  \mathbf{\Lambda}\right)  \right\Vert  &
=n_{j}\left(  \mathbf{\eta}\right)  ^{\frac{1}{2}}%
\end{align*}%
\[
c_{j}\left(  \mathbf{\eta}\right)  =\mathbf{n}_{j}^{\frac{1}{2}}\left(
\mathbf{\eta}\right)
\]


\subsection{The AGP Energy Functional}%

\begin{align*}
\mathcal{E}\left(  D^{N}\right)   &  =\frac{1}{\operatorname{Tr}\left\{
D^{N}\right\}  }\operatorname{Tr}\left\{  \left\{  \sum_{1\leq i,k\leq
2s}h_{1ik}a_{i}^{\dagger}a_{k}\right.  \right. \\
&  +\left.  \left.  \frac{1}{4}\sum_{1\leq i,j,k,l\leq2s}V_{2ijkl}%
a_{i}^{\dagger}a_{j}^{\dagger}a_{l}a_{k}\right\}  D^{N}\right\}
\end{align*}
%

\[
\mathcal{E}\left(  \mathbb{D}^{2}\right)  =\operatorname{Tr}\left\{
\mathbb{H}_{1}\mathbb{D}^{1}\right\}  +\frac{1}{4}\operatorname{Tr}\left\{
\mathbb{V}_{2}\mathbb{D}^{2}\right\}
\]%
\[
\left\langle jk||lm\right\rangle =\left(  \left.  jk\right\vert h_{2}%
lm\right)  -\left(  \left.  jk\right\vert h_{2}ml\right)  ,
\]%
\begin{align*}
\left\langle \varphi_{j}\varphi_{k}|||\varphi_{j}\varphi_{k}\right\rangle  &
=\left\langle jk||jk\right\rangle +\left\langle j\,k+s||j\,k+s\right\rangle \\
&  +\left\langle j+s\,k||j+s\,k\right\rangle +\left\langle
j+s\,k+s||j+s\,k+s\right\rangle
\end{align*}


\section{Optimization Theory\label{Optimiz}}%

\[
\left\langle \left.  \Upsilon_{j}\left(  \mathbf{\Lambda}\right)  \right\vert
\Upsilon_{k}\left(  \mathbf{\Lambda}\right)  \right\rangle =p_{j}\delta
_{jk};\quad p_{j}\geq0,\quad1\leq j\leq2s
\]%
\[
p_{j}\left(  \mathbf{\eta}\right)  =c_{0j}e^{\eta_{j}}\mathit{s}\left(
\mathbf{\eta}\right)  ^{-1}%
\]


\section{\bigskip AGP Lagrangian\label{Lagran}}

\section{\bigskip Generalized Janaks Theorem\label{General}}

\section{\bigskip Discussion\label{Discuss}}

The total energy can be expressed as
\begin{align*}
\mathcal{E}\left(  \mathbf{\varphi}\right)   &  =\frac{1}{2}\sum_{1\leq
p\leq2s}N\left(  \varphi_{p}\right)  \left\{  I\left(  \varphi_{p}\right)
+\left\langle \left.  \varphi_{p}\right\vert h_{1}\varphi_{p}\right\rangle
\right\} \\
&  =\frac{1}{2}%
%TCIMACRO{\tsum \limits_{1\leq p\leq2s}}%
%BeginExpansion
{\textstyle\sum\limits_{1\leq p\leq2s}}
%EndExpansion
N\left(  \varphi_{p}\right)  \kappa\left(  \varphi_{p}\right)
\end{align*}
where $\left\{  N\left(  \varphi_{p}\right)  \right\}  $ are the occupation
numbers of the orthogonal CGSO's $\left\{  \left\vert \varphi_{p}\right\rangle
\right\}  $ and the one electron quantities $\left\{  \kappa\left(
\varphi_{p}\right)  \right\}  $ are
\[
\kappa\left(  \varphi_{p}\right)  =I\left(  \varphi_{p}\right)  +\left\langle
\left.  \varphi_{p}\right\vert h_{1}\varphi_{p}\right\rangle
\]
in terms of the ionization potentials $\left\{  I\left(  \varphi_{p}\right)
\right\}  ,$ which are the energy differences between the normalized
$N-$electron AGP state and the normalized $N-1$ electron state formed by the
removal of the GSO $\left\vert \varphi_{p}\right\rangle ,$ and the one
electron Hamiltonian $h_{1}$
\[
h_{1}=\sum_{1\leq j,k\leq2s}h_{1jk}\left(  \mathbf{\varphi}\right)  \left\vert
\varphi_{j}\right\rangle \left\langle \varphi_{k}\right\vert ,
\]
where
\begin{align}
h_{1jk}\left(  \varphi\right)   &  =\left\langle \varphi_{j}\left\vert
h\right.  \varphi_{k}\right\rangle \label{h1}\\
&  =-\int\varphi_{j}^{\ast}\left(  \mathbf{r,\xi}\right)  \left\{
\sum_{_{1\leq l\leq N_{nuc}}}\frac{Z_{l}}{\left\Vert \mathbf{R}_{l}%
-\mathbf{r}\right\Vert }+\nabla_{\mathbf{r}}^{2}\right\}  \varphi_{k}\left(
\mathbf{r,\xi}\right)  d\mathbf{r}d\xi;1\leq j,k\leq2s.
\end{align}
Using the Fock operator the energy can be written as
\[
\mathcal{E}\left(  \mathbb{D}^{2},\mathbf{\varphi}\right)  =\frac{1}%
{2}\left\{  \overset{}{\operatorname{Tr}\left\{  F\left(  \mathbb{D}%
^{2},\mathbf{\varphi}\right)  \right\}  +\mathcal{E}_{1}\left(  \mathbb{D}%
^{2},\mathbf{\varphi}\right)  }\right\}  ,
\]
where
\[
\mathcal{E}_{1}\left(  \mathbb{D}^{2},\mathbf{\varphi}\right)
=\operatorname{Tr}\left\{  \mathbb{D}^{1}\left(  \mathbf{\varphi}\right)
\mathbb{H}_{1}\left(  \mathbf{\varphi}\right)  \right\}  .
\]
This gives
\[
\mathcal{E}\left(  \mathbb{D}^{2},\mathbf{\varphi}\right)  =\frac{1}{2}%
\sum_{1\leq p\leq2s}\left\{  \varepsilon\left(  \varphi_{p}\right)  +h\left(
\varphi_{p}\right)  N\left(  \varphi_{p}\right)  \right\}  ,
\]
For AGP states
\begin{align*}
\mathcal{E}\left(  \mathbf{\eta},\mathbf{\psi}\left(  \mathbf{\eta}\right)
\right)   &  =\frac{1}{2}\sum_{1\leq p\leq2s}\left\{  \varepsilon\left(
\mathbf{\eta,}\psi_{p}\left(  \mathbf{\eta}\right)  \right)  +h\left(
\psi_{p}\left(  \mathbf{\eta}\right)  \right)  N\left(  \psi_{p}\left(
\mathbf{\eta}\right)  \right)  \right\} \\
&  =\frac{1}{2}\sum_{1\leq p\leq2s}\left\{  \varepsilon\left(  \mathbf{\eta
,}\chi_{p}\left(  \mathbf{\eta}\right)  \right)  +h\left(  \psi_{p}\left(
\mathbf{\eta}\right)  \right)  N\left(  \psi_{p}\left(  \mathbf{\eta}\right)
\right)  \right\}
\end{align*}%
\[
\frac{\mathcal{E}\left(  \mathbf{\eta},\mathbf{\psi}\left(  \mathbf{\eta
}\right)  \right)  }{\partial N_{k}}=\frac{1}{2}\sum_{\substack{1\leq
p\leq2s\\1\leq l\leq s}}\left\{  \frac{\partial\varepsilon\left(
\mathbf{\eta,}\chi_{p}\left(  \mathbf{\eta}\right)  \right)  }{\partial
\eta_{l}}\frac{\partial\eta_{l}}{\partial N_{k}}\right\}  +\frac{1}{2}h\left(
\psi_{k}\left(  \mathbf{\eta}\right)  \right)
\]
where $\left\{  \varepsilon\left(  \mathbf{\eta,}\varphi_{p}\left(
\mathbf{\eta}\right)  \right)  \right\}  $ are diagonal elements of the Fock
operator in the basis $\left\{  \left\vert \varphi_{p}\right\rangle \right\}
$ and $\left\{  \varphi_{p}\right\}  $ are NGSO's. where $\left\{
\varepsilon\left(  \mathbf{\eta,}\psi_{p}\left(  \mathbf{\eta}\right)
\right)  \right\}  $ are diagonal elements of the Fock operator in the basis
$\left\{  \left\vert \psi_{p}\right\rangle \right\}  $ and $\left\{  \psi
_{p}\right\}  $ are CGSO's. where $\left\{  \varepsilon\left(  \mathbf{\eta
,}\chi_{p}\left(  \mathbf{\eta}\right)  \right)  \right\}  $ are diagonal
elements of the Fock operator in the basis $\left\{  \left\vert \chi
_{p}\right\rangle \right\}  $ and $\left\{  \chi_{p}\right\}  $ are the
eigenvectors of the Fock operator $F\left(  \mathbb{D}^{2},\mathbf{\varphi
}\right)  $. One should note that the GSO's that diagonalize the Fock operator
are \emph{not }in general the NGSO's. The eigenvalues of $F\left(
\mathbb{D}^{2},\mathbf{\varphi}\right)  $ are related to ionization
potentials
\[
\varepsilon\left(  \chi_{p}\right)  =N\left(  \chi_{p}\right)  I\left(
\chi_{p}\right)
\]
%

\[
\left\{  \left\vert \chi_{p}\right\rangle \right\}
\]


The energy of the AGP state is given by \emph{ }
\begin{align*}
&  \mathcal{E}\left(  \mathbf{n,\varphi}\right)  =\frac{\mathcal{H}\left(
\mathbf{n,\varphi}\right)  }{S_{\frac{N}{2}}\left(  \mathbf{n}\right)  }\\
&  =\frac{1}{S_{\frac{N}{2}}\left(  \mathbf{n}\right)  }\left\{  \sum_{1\leq
j\leq s}n_{j}S_{_{\frac{N}{2}-1}}\left(  \jmath\right)  \left\{  \left\langle
\varphi_{j}\left\vert h_{1}\varphi_{j}\right.  \right\rangle +\left\langle
\varphi_{j+s}\left\vert h_{1}\varphi_{j+s}\right.  \right\rangle +\left\langle
\varphi_{j}\varphi_{j+s}||\varphi_{j}\varphi_{j+s}\right\rangle \right\}
\right.  \\
&  +\sum_{1\leq j,k\leq s}n_{j}^{\frac{1}{2}}n_{k}^{\frac{1}{2}}S_{_{\frac
{N}{2}-1}}\left(  \jmath k\right)  \left\langle \varphi_{j}\varphi
_{j+s}||\varphi_{k}\varphi_{k+s}\right\rangle \left.  +\sum_{1\leq j<k\leq
s}n_{i}n_{j}S_{_{\frac{N}{2}-2}}\left(  \jmath k\right)  \left\langle
\varphi_{j}\varphi_{k}|||\varphi_{j}\varphi_{k}\right\rangle \right\}
\end{align*}%
\begin{align*}
&  \mathcal{E}\left(  \mathbf{n,\psi}\right)  =\frac{\mathcal{H}\left(
\mathbf{n,\psi}\right)  }{S_{\frac{N}{2}}\left(  \mathbf{n}\right)  }\\
&  =\frac{1}{S_{\frac{N}{2}}\left(  \mathbf{n}\right)  }\left\{  \sum_{1\leq
j\leq s}n_{j}S_{_{\frac{N}{2}-1}}\left(  \jmath\right)  \left\{  \left\langle
\psi_{j}\left\vert h_{1}\psi_{j}\right.  \right\rangle +\left\langle
\psi_{j+s}\left\vert h_{1}\psi_{j+s}\right.  \right\rangle +\left\langle
\psi_{j}\psi_{j+s}||\psi_{j}\psi_{j+s}\right\rangle \right\}  \right.  \\
&  +\sum_{1\leq j,k\leq s}n_{j}^{\frac{1}{2}}n_{k}^{\frac{1}{2}}S_{_{\frac
{N}{2}-1}}\left(  \jmath k\right)  \left\langle \psi_{j}\psi_{j+s}||\psi
_{k}\psi_{k+s}\right\rangle \left.  +\sum_{1\leq j<k\leq s}n_{i}%
n_{j}S_{_{\frac{N}{2}-2}}\left(  \jmath k\right)  \left\langle \psi_{j}%
\psi_{k}|||\psi_{j}\psi_{k}\right\rangle \right\}
\end{align*}
in terms of geminal occupation numbers and orthonormal CGSO's, while in terms
of the one electron group parameters it is
\begin{equation}
\mathcal{E}\left(  \mathbf{\eta},\mathbf{\lambda}\right)  =\frac
{\mathcal{H}\left(  \mathbf{\eta},\mathbf{\lambda}\right)  }{S_{\frac{N}{2}%
}\left(  \mathbf{\eta}\right)  }=\frac{\left\langle \left.  g\left(
\mathbf{\eta},\mathbf{\lambda}\right)  ^{\frac{N}{2}}\right\vert Hg\left(
\mathbf{\eta},\mathbf{\lambda}\right)  ^{\frac{N}{2}}\right\rangle }%
{S_{\frac{N}{2}}\left(  \mathbf{\eta}\right)  },\label{gp_funct}%
\end{equation}
where%
\begin{align*}
S_{\frac{N}{2}}\left(  \mathbf{\eta}\right)   &  =\left\langle \left.
k\left(  \mathbf{\eta}\right)  g_{0}^{\frac{N}{2}}\right\vert k\left(
\mathbf{\eta}\right)  g_{0}^{\frac{N}{2}}\right\rangle \\
&  =\mathit{s}\left(  \mathbf{\eta}\right)  ^{-\frac{N}{2}}\sum_{1\leq
j_{1}<\cdots<j_{\frac{N}{2}}\leq s}e^{2\left(  \eta_{j_{1}}+\cdots+\eta
_{j_{N}}\right)  }n_{0j_{1}}\cdots n_{0j_{N}}%
\end{align*}%
\begin{equation}
n_{j}\left(  \mathbf{\eta}\right)  =e^{2\eta_{j}}n_{0j}\mathit{s}\left(
\mathbf{\eta}\right)  ^{-\frac{1}{2}}%
\end{equation}%
\begin{align*}
&  \mathcal{E}\left(  \mathbf{\Upsilon}\left(  \mathbf{\Lambda}\right)
\right)  =\mathcal{E}\left(  \mathbf{\Lambda}\right)  =\frac{\mathcal{H}%
\left(  \mathbf{\Upsilon}\left(  \mathbf{\Lambda}\right)  \right)  }%
{S_{\frac{N}{2}}\left(  \mathbf{\eta}\right)  }\\
&  =\frac{1}{S_{\frac{N}{2}}\left(  \mathbf{\eta}\right)  }\left\{
\sum_{1\leq j\leq s}n_{j}\left(  \mathbf{\eta}\right)  ^{\frac{1}{2}}%
S_{\frac{N}{2}-1}\left(  \jmath\right)  \left\{  \left\langle \Upsilon
_{j}\left(  \mathbf{\Lambda}\right)  \left\vert h_{1}\Upsilon_{j}\left(
\mathbf{\Lambda}\right)  \right.  \right\rangle \right.  \right.  \\
&  \left.  +\quad\left\langle \Upsilon_{j+s}\left(  \mathbf{\Lambda}\right)
\left\vert h_{1}\Upsilon_{j+s}\left(  \mathbf{\Lambda}\right)  \right.
\right\rangle +\,\left\langle \Upsilon_{j}\left(  \mathbf{\Lambda}\right)
\Upsilon_{j+s}\left(  \mathbf{\Lambda}\right)  ||\Upsilon_{j}\left(
\mathbf{\Lambda}\right)  \Upsilon_{j+s}\left(  \mathbf{\Lambda}\right)
\right\rangle \right\}  \\
&  +\sum_{1\leq j,k\leq s}S_{\frac{N}{2}-1}\left(  \jmath k\right)
\left\langle \Upsilon_{j}\left(  \mathbf{\Lambda}\right)  \Upsilon
_{j+s}\left(  \mathbf{\Lambda}\right)  ||\Upsilon_{k}\left(  \mathbf{\Lambda
}\right)  \Upsilon_{k+s}\left(  \mathbf{\Lambda}\right)  \right\rangle \\
&  \left.  +\sum_{1\leq j<k\leq s}n_{j}\left(  \mathbf{\eta}\right)
^{\frac{1}{2}}n_{k}\left(  \mathbf{\eta}\right)  ^{\frac{1}{2}}S_{\frac{N}%
{2}-2}\left(  \jmath k\right)  \left\langle \Upsilon_{j}\left(
\mathbf{\Lambda}\right)  \Upsilon_{k}\left(  \mathbf{\Lambda}\right)
|||\Upsilon_{j}\left(  \mathbf{\Lambda}\right)  \Upsilon_{k}\left(
\mathbf{\Lambda}\right)  \right\rangle \right\}
\end{align*}
The constraints on the variation are
\[
\left\langle \left.  \Upsilon_{j}\left(  \mathbf{\Lambda}\right)  \right\vert
\Upsilon_{k}\left(  \mathbf{\Lambda}\right)  \right\rangle =c_{0j}e^{\eta_{j}%
}\mathit{s}\left(  \mathbf{\eta}\right)  ^{-1}\delta_{jk};\quad-\infty\leq
\eta_{j}\leq\infty,\quad1\leq j\leq2s.
\]
for fixed $\mathbf{\eta}$ the canonical coefficients $\mathbf{c}\left(
\mathbf{\eta}\right)  $ are fixed. Thus one can optimize $\mathcal{E}\left(
\mathbf{\Lambda}\right)  $ by optimizing $\mathbf{\theta}$ and $\mathbf{\alpha
}$ for fixed coefficients, then relaxing the normalization of $\left\{
\Upsilon_{j}\left(  \mathbf{\Lambda}\right)  \right\}  $ and again optimizing
$\mathbf{\theta}$ and $\mathbf{\alpha}$ etc. The normalization of $\left\{
\Upsilon_{j}\left(  \mathbf{\Lambda}\right)  \right\}  $ is in $1-1$
correspondence with $\mathbf{\eta}$. One can also consider this as an
unconstrained variation by just varying $\mathbf{\eta\in}\left[
-\mathbf{\infty,\infty}\right]  $ and $\mathbf{\theta}$ and $\mathbf{\alpha}$
as described below
\[
0\leq\theta_{j}<2\pi;\quad1\leq j\leq s
\]
as the energy functional is periodic in $\mathbf{\theta}$ this is not a
constriant.
\begin{align*}
\alpha_{jk} &  =\alpha_{kj}^{\ast};\quad1\leq j,k\leq2s\\
\alpha_{jk} &  =0;\quad j=k\pm s;\,j=k
\end{align*}
This constraint can be transformed away by setting
\[
\alpha_{jk}=x_{jk}+iy_{jk};\,1\leq j\neq k,k\pm s\leq2s
\]
so that
\begin{equation}
\alpha_{kj}^{\ast}=x_{kj}-iy_{kj}%
\end{equation}
leading to
\begin{align}
x_{kj} &  =x_{jk}\\
y_{jk} &  =-y_{kj}%
\end{align}
so that the constrained complex variables $\left\{  \alpha_{jk};1\leq
j,k\leq2s;\,j\neq k\pm s;j\neq k\right\}  $ can be replaced by unconstrained
real ones $\left\{  \left.  x_{jk},y_{kj}\right\vert 1\leq j<k\leq2s;\,j\neq
k\pm s\right\}  $ so that
\[
\mathbf{\alpha=\alpha}\left(  \mathbf{x},\mathbf{y}\right)
\]
Thus we can express the optimization problem in two ways
\[
\mathcal{E}\left(  \mathbf{\Upsilon}_{GS}\right)  =\min_{\mathbf{\eta}%
}\left\{  \min_{\mathbf{\Upsilon\in\Delta}\left(  \mathbf{\eta}\right)
}\mathcal{E}\left(  \mathbf{\Upsilon}\right)  \right\}
\]
where
\[
\mathbf{\Delta}\left(  \mathbf{\eta}\right)  =\left\{  \left\vert \Upsilon
_{j}\right\rangle \left\vert \overset{}{\left\langle \left.  \Upsilon
_{j}\right\vert \Upsilon_{k}\right\rangle =e^{\eta_{j}}\mathit{s}\left(
\mathbf{\eta}\right)  ^{-\frac{1}{2}}c_{0j}\delta_{jk};\left\vert \Upsilon
_{j}\right\rangle \in\mathcal{H}^{1}}\right.  \right\}
\]%
\[
\mathbf{\eta\in}\mathbb{R}^{s}%
\]%
\[
\mathcal{E}\left(  \mathbf{\Upsilon}_{GS}\right)  =\min_{\mathbf{c}}\left\{
\min_{\mathbf{\Upsilon\in\Delta}\left(  \mathbf{c}\right)  }\mathcal{E}\left(
\mathbf{\Upsilon}\right)  \right\}
\]
*******%
\[
\mathbf{\Delta}\left(  \mathbf{c}\right)  =\left\{  \left\vert \Upsilon
_{j}\right\rangle \left\vert \overset{}{\left\langle \left.  \Upsilon
_{j}\right\vert \Upsilon_{k}\right\rangle =c_{j}\delta_{jk};\left\vert
\Upsilon_{j}\right\rangle \in\mathcal{H}^{1}}\right.  \right\}
\]%
\[
\mathbf{c\geq0,\quad c\in}\mathbb{R}^{s}%
\]
and
\[
\mathcal{E}\left(  \mathbf{\Lambda}_{GS}\right)  =\min_{\mathbf{\Lambda\in
}\mathbb{R}^{s}\times\mathbb{R}^{u}\times\mathbb{R}^{s}}\mathcal{E}\left(
\mathbf{\Lambda}\right)
\]
where
\[
\mathbf{\Lambda}=\left(  \mathbf{\eta,\alpha}\left(  \mathbf{x,y}\right)
,\mathbf{\theta}\right)  ;\quad\mathbf{\eta\in}\mathbb{R}^{s},\mathbf{x,y\in
}\mathbb{R}^{u},\mathbf{\theta\in}\mathbb{R}^{s}%
\]%
\begin{equation}
u=\left(  2s+1\right)  -3s
\end{equation}
One can define
\[
\Omega=\left\{  \left.  \mathbf{\Upsilon}_{GS}\left(  \mathbf{p}\right)
=\min_{\mathbf{\Upsilon\in\Delta}\left(  \mathbf{p}\right)  }\mathcal{E}%
\left(  \mathbf{\Upsilon}\right)  \right\vert \mathbf{p\in}\mathbb{R}_{+}%
^{s}\right\}
\]
so that the ground energy is given by
\[
\mathcal{E}\left(  \mathbf{\Upsilon}_{GS}\left(  \mathbf{p}_{\#}\right)
\right)  =\min_{\mathbf{p\in}\mathbb{R}_{+}^{s}}\mathcal{E}\left(
\mathbf{\Upsilon}_{GS}\left(  \mathbf{p}\right)  \right)
\]
It is possible that for a given $\mathbf{p}$ the minimum is degenerate so one
obtains solutions $\left\{  \mathbf{\Upsilon}_{GS1}\left(  \mathbf{p}\right)
,\cdots,\mathbf{\Upsilon}_{GSm}\left(  \mathbf{p}\right)  \right\}  ,$ which
also must be included in the set $\Omega$. One can define a Lagrangian \#\#\#
\[
\mathcal{L}\left(  \mathbf{\Upsilon}\right)  =\mathcal{E}\left(
\mathbf{\Upsilon}\right)  -\sum_{1\leq j,k\leq2s}\varsigma_{jk}\left(
\left\langle \left.  \Upsilon_{j}\right\vert \Upsilon_{k}\right\rangle
-p_{j}\delta_{jk}\right)
\]
associated with the constrained variation (orthogonal with a fixed
normalization) of the energy functional $\mathcal{E}\left(  \mathbf{\Upsilon
}\right)  ,$ where
\[
\mathcal{L}:\overset{s}{\overbrace{\mathcal{H}^{1}\times\cdots\times
\mathcal{H}^{1}}}\times\,\mathbb{R\times R}\rightarrow\mathbb{R}%
\]
The stationarity conditions are
\[
\frac{\delta\mathcal{L}\left(  \mathbf{\Upsilon}\right)  }{\delta\Upsilon
_{j}^{\ast}}=\frac{\delta\mathcal{E}\left(  \mathbf{\Upsilon}\right)  }%
{\delta\Upsilon_{j}^{\ast}}-\sum_{1\leq k\leq2s}\varsigma_{jk}\left\vert
\Upsilon_{k}\right\rangle =0
\]
cc%
\[
\left\langle \left.  \Upsilon_{j}\left(  \mathbf{c}\right)  \right\vert
\Upsilon_{k}\left(  \mathbf{c}\right)  \right\rangle -c_{j}\delta_{jk}=0
\]%
\[
p_{j}=c_{0j}e^{\eta_{j}}\mathit{s}\left(  \mathbf{\eta}\right)  ^{-\frac{1}%
{2}}%
\]
The solutions $\left\{  \mathbf{\Upsilon}\left(  \mathbf{c}\right)  \right\}
$ and $\left\{  \mathbf{\varsigma}\left(  \mathbf{c}\right)  \right\}  $
obviously depend on $\mathbf{c}$. The sensitivity to the normalization
constraints is%
\begin{equation}
\frac{\partial\mathcal{L}\left(  \mathbf{\Upsilon}\right)  }{\partial p_{j}%
}=\varsigma_{jj}%
\end{equation}%
\begin{align*}
\frac{\partial\mathcal{L}\left(  \mathbf{\Upsilon}\right)  }{\partial\eta_{j}}
&  =\sum_{1\leq k\leq2s}\varsigma_{kk}c_{0k}\frac{\partial}{\partial\eta_{j}%
}\left\{  e^{\eta_{k}}\mathit{s}\left(  \mathbf{\eta}\right)  ^{-\frac{1}{2}%
}\right\}  \\
&  =\sum_{1\leq k\leq2s}\varsigma_{kk}c_{0k}\left\{  e^{\eta_{k}}\delta
_{jk}\mathit{s}\left(  \mathbf{\eta}\right)  ^{-\frac{1}{2}}+e^{\eta_{k}}%
\frac{\partial}{\partial\eta_{j}}\left\{  \mathit{s}\left(  \mathbf{\eta
}\right)  ^{-\frac{1}{2}}\right\}  \right\}  \\
&  =\varsigma_{jj}c_{0j}e^{\eta_{j}}\mathit{s}\left(  \mathbf{\eta}\right)
^{-\frac{1}{2}}-\frac{1}{2}\sum_{1\leq k\leq2s}\varsigma_{kk}c_{0k}e^{\eta
_{k}}\mathit{s}\left(  \mathbf{\eta}\right)  ^{-\frac{3}{2}}\frac
{\partial\mathit{s}\left(  \mathbf{\eta}\right)  }{\partial\eta_{j}}\\
&  =\frac{1}{2}\mathit{s}\left(  \mathbf{\eta}\right)  ^{-\frac{1}{2}}\left\{
2\varsigma_{jj}c_{0j}e^{\eta_{j}}-\left(  \sum_{1\leq k\leq2s}\varsigma
_{kk}c_{0k}e^{\eta_{k}}\right)  \mathit{s}\left(  \mathbf{\eta}\right)
^{-1}\frac{\partial\mathit{s}\left(  \mathbf{\eta}\right)  }{\partial\eta_{j}%
}\right\}  \\
&  \frac{1}{2}\mathit{s}\left(  \mathbf{\eta}\right)  ^{-\frac{1}{2}}\left\{
2\varsigma_{jj}c_{0j}e^{\eta_{j}}-\left(  \sum_{1\leq k\leq2s}\varsigma
_{kk}c_{0k}e^{\eta_{k}}\right)  \mathit{s}\left(  \mathbf{\eta}\right)
^{-1}n_{0j}e^{2\eta_{j}}\right\}
\end{align*}%
\[
\frac{\partial\mathcal{L}\left(  \mathbf{\Upsilon}\right)  }{\partial\eta_{j}%
}=\sum_{1\leq k\leq2s}\varsigma_{kk}c_{0k}\frac{\partial}{\partial\eta_{j}%
}\left\{  e^{\eta_{k}}\mathit{s}\left(  \mathbf{\eta}\right)  ^{-\frac{1}{2}%
}\right\}
\]%
\[
\mathit{s}\left(  \mathbf{\eta}\right)  =\sum_{1\leq k\leq2s}n_{k}\left(
\mathbf{\eta}\right)  =\sum_{1\leq k\leq2s}n_{0k}e^{2\eta_{k}}%
\]%
\begin{equation}
\frac{\partial\mathit{s}\left(  \mathbf{\eta}\right)  }{\partial\eta_{j}%
}=n_{0j}e^{2\eta_{j}}%
\end{equation}%
\begin{equation}
2\mathit{s}\left(  \mathbf{\eta}\right)  =\frac{\partial\mathit{s}\left(
\mathbf{\eta}\right)  }{\partial\eta_{j}}\Rightarrow2\sum_{1\leq k\leq
2s}n_{0k}e^{2\eta_{k}}=n_{0j}e^{2\eta_{j}}\forall j
\end{equation}%
\[
\frac{\partial\mathcal{L}\left(  \mathbf{\Upsilon}\right)  }{\partial\eta_{j}%
}=\sum_{1\leq k\leq2s}\varsigma_{kk}c_{0k}\frac{\partial}{\partial\eta_{j}%
}\left\{  e^{\eta_{k}}\mathit{s}\left(  \mathbf{\eta}\right)  ^{-\frac{1}{2}%
}\right\}
\]%
\[
\frac{\partial\mathcal{L}\left(  \mathbf{\Upsilon}\right)  }{\partial c_{j}%
}=\varsigma_{kk}\left(  \mathbf{c}\right)
\]


Restricting the functionals $\mathcal{L}$ and $\mathcal{E}$

$\qquad$to the set $\Omega$ one obtains%
\begin{equation}
\mathcal{L}\left(  \mathbf{\Upsilon}_{GS}\left(  \mathbf{c}\right)  \right)
=\mathcal{E}\left(  \mathbf{\Upsilon}_{GS}\left(  \mathbf{c}\right)  \right)
\end{equation}%
\[
\frac{\partial\mathcal{L}\left(  \mathbf{\Upsilon}_{GS}\left(  \mathbf{\eta
}\right)  \right)  }{\partial\eta_{j}}=\frac{\partial\mathcal{E}\left(
\mathbf{\Upsilon}_{GS}\left(  \mathbf{\eta}\right)  \right)  }{\partial
\eta_{j}}%
\]%
\[
\frac{\partial\mathcal{L}\left(  \mathbf{\Upsilon}_{GS}\left(  \mathbf{c}%
\right)  \right)  }{\partial c_{j}}=\frac{\partial\mathcal{E}\left(
\mathbf{\Upsilon}_{GS}\left(  \mathbf{c}\right)  \right)  }{\partial c_{j}}%
\]
as
\begin{align*}
\mathcal{L}\left(  \mathbf{\Upsilon}_{GS}\left(  \mathbf{\eta}\right)
\right)   &  =\mathcal{E}\left(  \mathbf{\Upsilon}_{GS}\left(  \mathbf{\eta
}\right)  \right) \\
&  -\sum_{1\leq j,k\leq2s}\varsigma_{kj}\left(  \left\langle \left.
\Upsilon_{GSj}\left(  \mathbf{\eta}\right)  \right\vert \Upsilon_{GSk}\left(
\mathbf{\eta}\right)  \right\rangle -c_{0j}e^{\eta_{j}}\delta_{jk}\right)
\end{align*}
and
\[
\sum_{1\leq j,k\leq2s}\varsigma_{kj}\left(  \left\langle \left.
\Upsilon_{GSj}\left(  \mathbf{\eta}\right)  \right\vert \Upsilon_{GSk}\left(
\mathbf{\eta}\right)  \right\rangle -c_{0j}e^{\eta_{j}}\delta_{jk}\right)  =0
\]
Thus in particular
\[
\frac{\partial\mathcal{L}\left(  \mathbf{\Upsilon}_{GS}\left(  \mathbf{\eta
}\right)  \right)  }{\partial\eta_{j}}=\frac{\partial\mathcal{E}\left(
\mathbf{\Upsilon}_{GS}\left(  \mathbf{\eta}\right)  \right)  }{\partial
\eta_{j}}=\varsigma_{jj}\left\langle \left.  \Upsilon_{0j}\right\vert
\Upsilon_{0j}\right\rangle
\]%
\[
\frac{\partial\mathcal{L}\left(  \mathbf{\Upsilon}\left(  \mathbf{c}\right)
\right)  }{\partial c_{j}}=\frac{\partial\mathcal{E}\left(  \mathbf{\Upsilon
}\left(  \mathbf{c}\right)  \right)  }{\partial c_{j}}=\varsigma_{jj}\left(
\mathbf{c}\right)
\]
As $\mathbf{N}\leftrightarrow\mathbf{n}$ $\mathbf{\leftrightarrow c}$ one
obtains
\begin{align*}
\frac{\partial\mathcal{E}\left(  \mathbf{\Upsilon}\right)  }{\partial N_{j}}
&  =\sum_{1\leq k\leq2s}\frac{\partial\mathcal{E}\left(  \mathbf{\Upsilon
}\right)  }{\partial n_{k}}\frac{\partial n_{k}}{\partial N_{j}}\\
&  =\sum_{1\leq k\leq2s}\frac{\partial\mathcal{E}\left(  \mathbf{\Upsilon
}\right)  }{\partial c_{k}}\frac{\partial c_{k}}{\partial n_{k}}\frac{\partial
n_{k}}{\partial N_{j}}\\
&  =\sum_{1\leq k\leq2s}\frac{\partial\mathcal{E}\left(  \mathbf{\Upsilon
}\right)  }{\partial c_{k}}\frac{1}{2c_{k}}\frac{\partial n_{k}}{\partial
N_{j}}\\
&  =\sum_{1\leq k\leq2s}\frac{\partial\mathcal{E}\left(  \mathbf{\Upsilon
}\right)  }{\partial c_{k}}\frac{1}{2c_{k}}\left(  \frac{\partial\mathbf{N}%
}{\partial\mathbf{n}}\right)  _{kj}^{-1}%
\end{align*}
JJ%
\begin{align*}
\frac{\partial\mathcal{E}\left(  \mathbf{\Upsilon}\right)  }{\partial N_{j}}
&  =\sum_{1\leq k\leq2s}\frac{\partial\mathcal{E}\left(  \mathbf{\Upsilon
}\right)  }{\partial c_{k}}\frac{\partial c_{k}}{\partial n_{k}}\frac{\partial
n_{k}}{\partial N_{j}}\\
&  =\sum_{1\leq k\leq2s}\frac{\partial\mathcal{E}\left(  \mathbf{\Upsilon
}\right)  }{\partial c_{k}}\frac{1}{2c_{k}}\left(  \frac{\partial\mathbf{N}%
}{\partial\mathbf{n}}\right)  _{kj}^{-1}%
\end{align*}
i.e.%
\[
\frac{\partial\mathcal{E}\left(  \mathbf{\Upsilon}\right)  }{\partial
\mathbf{N}}=\frac{1}{2}\frac{\partial\mathcal{E}\left(  \mathbf{\Upsilon
}\right)  }{\partial\mathbf{c}}\mathbf{c}^{-1}\left(  \frac{\partial
\mathbf{N}}{\partial\mathbf{\eta}}\right)  ^{-1}=\frac{1}{2}\frac
{\partial\mathcal{E}\left(  \mathbf{\Upsilon}\right)  }{\partial\mathbf{c}%
}\left(  \frac{\partial\mathbf{N}}{\partial\mathbf{\eta}}\mathbf{c}\right)
^{-1}%
\]%
\[
\frac{\partial\mathcal{E}\left(  \mathbf{\Upsilon}\right)  }{\partial
\mathbf{N}}=\frac{1}{2}\frac{\partial\mathcal{E}\left(  \mathbf{\Upsilon
}\right)  }{\partial\mathbf{c}}\left(  \frac{\partial\mathbf{N}}%
{\partial\mathbf{\eta}}\mathbf{c}\right)  ^{-1}%
\]%
\begin{align*}
&  \quad\quad\left.  \Updownarrow\right. \\
\frac{\partial\mathcal{E}\left(  \mathbf{\Upsilon}\right)  }{\partial
\mathbf{c}}  &  =2\frac{\partial\mathcal{E}\left(  \mathbf{\Upsilon}\right)
}{\partial\mathbf{N}}\left(  \frac{\partial\mathbf{N}}{\partial\mathbf{\eta}%
}\mathbf{c}\right)
\end{align*}%
\[
\frac{\partial\mathcal{E}\left(  \mathbf{\Upsilon}\left(  \mathbf{c}%
_{GS}\right)  \right)  }{\partial\mathbf{c}}=\mathbf{0}%
\]%
\[
\mathbf{\varsigma}\left(  \mathbf{c}_{GS}\right)  =\mathbf{0}%
\]
The occupation numbers of the AGP\ state are
\[
N_{j}=\mathcal{N}_{j}\left(  \mathbf{n}\right)  =n_{j}\frac{S_{\frac{N}{2}%
-1}\left(  \jmath\right)  }{S_{\frac{N}{2}}}=n_{j}\frac{\partial\ln
S_{\frac{N}{2}}}{\partial n_{j}};\quad1\leq j\leq s,
\]%
\[
\mathbf{N}=\mathcal{N}\left(  \mathbf{n}\right)
\]%
\begin{align*}
&  \mathcal{N}_{j}\left(  \mathbf{n}\left(  \mathbf{\eta}\right)  \right) \\
&  \,\,\,\,\,\,\,\,\,\,\left.  \shortparallel\right.
\end{align*}%
\[
\frac{e^{2\eta_{j}}n_{0j}%
%TCIMACRO{\dsum \limits_{_{\substack{1\leq j_{1}<\cdots<j_{\frac{N}{2}-1}\leq
%s\\j\notin\left\{  j_{1},\cdots,j_{\frac{N}{2}-1}\right\}  }}}}%
%BeginExpansion
{\displaystyle\sum\limits_{_{\substack{1\leq j_{1}<\cdots<j_{\frac{N}{2}%
-1}\leq s\\j\notin\left\{  j_{1},\cdots,j_{\frac{N}{2}-1}\right\}  }}}}
%EndExpansion
e^{2\left(  \eta_{j_{1}}+\cdots+\eta_{j_{\frac{N}{2}-1}}\right)  }n_{0j_{1}%
}\cdots n_{0j_{\frac{N}{2}}}}{%
%TCIMACRO{\dsum \limits_{1\leq j_{1}<\cdots<j_{\frac{N}{2}}\leq s}}%
%BeginExpansion
{\displaystyle\sum\limits_{1\leq j_{1}<\cdots<j_{\frac{N}{2}}\leq s}}
%EndExpansion
e^{2\left(  \eta_{j_{1}}+\cdots+\eta_{j_{\frac{_{N}}{2}}}\right)  }n_{0j_{1}%
}\cdots n_{0j_{\frac{N}{2}}}}%
\]%
\begin{align*}
&  S_{\frac{N}{2}}\left(  \mathbf{n}\left(  \mathbf{\eta}\right)  \right) \\
&  \,\,\,\,\,\,\,\,\,\,\left.  \shortparallel\right.
\end{align*}%
\[
\mathit{s}\left(  \mathbf{\eta}\right)  ^{-\frac{N}{2}}\sum_{1\leq
j_{1}<\cdots<j_{\frac{N}{2}}\leq s}e^{2\left(  \eta_{j_{1}}+\cdots
+\eta_{j\frac{_{N}}{2}}\right)  }n_{0j_{1}}\cdots n_{0j_{\frac{N}{2}}}%
\]%
\begin{align}
&  \mathit{s}\left(  \mathbf{\eta}\right)  ^{-1}\\
n_{j}\left(  \mathbf{\eta}\right)   &  =e^{2\eta_{j}}n_{0j}\mathit{s}\left(
\mathbf{\eta}\right)  ^{-1}\\
&  =\mathit{s}\left(  \mathbf{\eta}\right)  ^{-\frac{N}{2}}\sum_{1\leq
j_{1}<\cdots<j_{\frac{N}{2}}\leq s}e^{2\left(  \eta_{j_{1}}+\cdots
+\eta_{j_{\frac{_{N}}{2}}}\right)  }n_{0j_{1}}\cdots n_{0j_{\frac{N}{2}}}%
\end{align}%
\[
n_{j}\left(  \mathbf{\eta}\right)  =e^{2\eta_{j}}n_{0j}\mathit{s}\left(
\mathbf{\eta}\right)  ^{-1}%
\]
The occupation number variation nonunitary subgroup $G_{num}$ is given by
\begin{equation}
G_{num}=\left\{  \left.  k\left(  \mathbf{\eta}\right)  =\exp\left\{
\sum_{1\leq j\leq s}\eta_{j}\Omega_{j}\right\}  \right\vert \eta_{j}%
\in\mathbb{R}\right\}  , \label{occup}%
\end{equation}
where
\begin{equation}
\Omega_{j}=\left(  a_{j}^{\dagger}a_{j}+a_{j+s}^{\dagger}a_{j+s}\right)
;\quad1\leq j\leq s.
\end{equation}
The manifold of $2N$-electron AGP\ states $\mathcal{M}$ can be expressed as
\begin{equation}
\mathcal{M}\mathcal{=}\left\{  \left.  \overset{}{k\left(  \mathbf{\eta
}\right)  u\left(  \mathbf{\lambda}\right)  }\left\vert g_{0}^{N}\right\rangle
\right\vert \right.  \left.  \mathbf{\lambda}=\mathbf{\lambda}^{\dagger
}\right\}  .
\end{equation}
Explicitly the action on a reference geminal is given by
\begin{equation}
\left\vert \tilde{g}\left(  \mathbf{\eta},\mathbf{\lambda}\right)
\right\rangle =k\left(  \mathbf{\eta}\right)  u\left(  \mathbf{\lambda
}\right)  \left\vert g_{0}\right\rangle =\sum_{1\leq l\leq s}e^{\eta_{l}%
}c_{0l}\left\vert \varphi_{l}\left(  \mathbf{\lambda}\right)  \varphi
_{l+s}\left(  \mathbf{\lambda}\right)  \right\rangle ,
\end{equation}
where
\begin{equation}
\left\vert \varphi_{l}\left(  \mathbf{\lambda}\right)  \right\rangle
=\exp\left\{  i\sum_{1\leq j,k\leq2s}\lambda_{jk}a_{j}^{\dagger}a_{k}\right\}
\left\vert \varphi_{0l}\right\rangle .
\end{equation}
One can define a normalized geminal by
\begin{equation}
\left\vert g\left(  \mathbf{\eta},\mathbf{\lambda}\right)  \right\rangle
=\mathit{s}\left(  \mathbf{\eta}\right)  ^{-\frac{1}{2}}\left\vert \tilde
{g}\left(  \mathbf{\eta},\mathbf{\lambda}\right)  \right\rangle ,
\end{equation}
where
\begin{equation}
\mathit{s}\left(  \mathbf{\eta}\right)  =\left\langle \left.  \tilde{g}\left(
\mathbf{\eta},\mathbf{\lambda}\right)  \right\vert \tilde{g}\left(
\mathbf{\eta},\mathbf{\lambda}\right)  \right\rangle =\sum_{1\leq l\leq
s}e^{2\eta_{l}}c_{0l}^{2}=\sum_{1\leq l\leq s}e^{2\eta_{l}}n_{0l}%
\end{equation}
giving
\begin{align}
\left\vert g\left(  \mathbf{\eta},\mathbf{\lambda}\right)  \right\rangle  &
=\mathit{s}\left(  \mathbf{\eta}\right)  ^{-\frac{1}{2}}\sum_{1\leq l\leq
s}e^{\eta_{l}}c_{0l}\left\vert \varphi_{l}\left(  \mathbf{\lambda}\right)
\varphi_{l+s}\left(  \mathbf{\lambda}\right)  \right\rangle \\
&  =\sum_{1\leq l\leq s}e^{\eta_{l}}c_{0l}\mathit{s}\left(  \mathbf{\eta
}\right)  ^{-\frac{1}{2}}\left\vert \varphi_{l}\left(  \mathbf{\lambda
}\right)  \varphi_{l+s}\left(  \mathbf{\lambda}\right)  \right\rangle
=\sum_{1\leq l\leq s}c_{l}\left(  \mathbf{\eta}\right)  \left\vert \varphi
_{l}\left(  \mathbf{\lambda}\right)  \varphi_{l+s}\left(  \mathbf{\lambda
}\right)  \right\rangle ,
\end{align}
where
\[
c_{l}\left(  \mathbf{\eta}\right)  =e^{\eta_{l}}c_{0l}\mathit{s}\left(
\mathbf{\eta}\right)  ^{-\frac{1}{2}}%
\]
leading to
\[
n_{l}\left(  \mathbf{\eta}\right)  =e^{2\eta_{l}}n_{0l}\mathit{s}\left(
\mathbf{\eta}\right)  ^{-1}.
\]
This gives
\begin{equation}
\sum_{1\leq l\leq s}n_{l}\left(  \mathbf{\eta}\right)  =\mathit{s}\left(
\mathbf{\eta}\right)  ^{-1}\sum_{1\leq l\leq s}e^{2\eta_{l}}n_{0l}=\left(
\sum_{1\leq l\leq s}e^{2\eta_{l}}n_{0l}\right)  ^{-1}\sum_{1\leq l\leq
s}e^{2\eta_{l}}n_{0l}=1
\end{equation}
which shows that $\left\{  n_{l}\left(  \mathbf{\eta}\right)  \right\}  $
satisfy the geminal occupation number positivity and normalization
constraints. We can define%
\begin{equation}
\left\vert \varphi_{l}\left(  \mathbf{\eta,\lambda}\right)  \right\rangle
=\exp\eta_{l}\exp\left\{  i\sum_{1\leq j,k\leq2s}\lambda_{jk}a_{j}^{\dagger
}a_{k}\right\}  \left\vert \varphi_{0l}\right\rangle ,
\end{equation}
so that
\begin{equation}
\left\langle \left.  \varphi_{l}\left(  \mathbf{\eta,\lambda}\right)
\right\vert \varphi_{l}\left(  \mathbf{\eta,\lambda}\right)  \right\rangle
=e^{2\eta_{l}}\left\langle \left.  \varphi_{0l}\right\vert \varphi
_{0l}\right\rangle =e^{2\eta_{l}}%
\end{equation}
and
\begin{equation}
\left\vert g\left(  \mathbf{\eta},\mathbf{\lambda}\right)  \right\rangle
=\mathit{s}\left(  \mathbf{\eta}\right)  ^{-\frac{1}{2}}\sum_{1\leq l\leq
s}c_{0l}\left\vert \varphi_{l}\left(  \mathbf{\eta,\lambda}\right)
\varphi_{l+s}\left(  \mathbf{\eta,\lambda}\right)  \right\rangle .
\end{equation}


We can express the geminal occupation numbers (and hence the occupation
numbers of the FORDM) and the normalization of the CGSO's in terms of
$\mathbf{\eta}$
\begin{align}
n_{l}\left(  \mathbf{\eta}\right)   &  =e^{2\eta_{l}}n_{0l}\mathit{s}\left(
\mathbf{\eta}\right)  ^{-1}\nonumber\\
s_{l}  &  =\left\langle \left.  \varphi_{l}\left(  \mathbf{\eta,\lambda
}\right)  \right\vert \varphi_{l}\left(  \mathbf{\eta,\lambda}\right)
\right\rangle =e^{2\eta_{l}}%
\end{align}
giving
\begin{equation}
\eta_{l}=\frac{1}{2}\ln\left\langle \left.  \varphi_{l}\left(  \mathbf{\eta
,\lambda}\right)  \right\vert \varphi_{l}\left(  \mathbf{\eta,\lambda}\right)
\right\rangle \equiv\frac{1}{2}\ln\left\langle \left.  \varphi_{l}\left(
\mathbf{\lambda}\right)  \right\vert \varphi_{l}\left(  \mathbf{\lambda
}\right)  \right\rangle =\frac{1}{2}\ln s_{l}%
\end{equation}
Hence%
\begin{equation}
\mathcal{E}\left(  \mathbf{n}_{\min}\mathbf{,\varphi}_{\min}\right)
=\mathcal{E}\left(  \mathbf{\eta}_{\min}\mathbf{,\lambda}_{\min}\right)
=\mathcal{E}\left(  \mathbf{\varphi}_{\min}\right)
\end{equation}
where
\begin{align}
\mathcal{E}\left(  \mathbf{n}_{\min}\mathbf{,\varphi}_{\min}\right)   &
=\min_{\mathbf{n\in\mathcal{S},\varphi\in}\mathcal{O}\left(  \mathbf{1}%
_{s}\right)  }\mathcal{E}\left(  \mathbf{n,\varphi}\right) \\
\mathcal{E}\left(  \mathbf{\eta}_{\min}\mathbf{,\lambda}_{\min}\right)   &
=\min_{\mathbf{\eta,\lambda}}\mathcal{E}\left(  \mathbf{\eta,\lambda}\right)
\\
\mathcal{E}\left(  \mathbf{\varphi}_{\min}\right)   &  =\min_{\mathbf{s\in
}\mathbb{R}_{+}^{s}}\min_{\mathbf{\varphi\in}\mathcal{O}\left(  \mathbf{s}%
\right)  }\mathcal{E}\left(  \mathbf{\varphi}\right)
\end{align}
and
\begin{align}
\mathcal{S}  &  =\left\{  \left.  \mathbf{n}\right\vert n_{j}\geq0,1\leq j\leq
s;\sum_{1\leq j\leq s}n_{j}=1\right\} \\
\mathcal{O}\left(  \mathbf{s}\right)   &  =\left\{  \left.  \mathbf{\varphi
}\right\vert \left\langle \left.  \varphi_{j}\right\vert \varphi
_{k}\right\rangle =s_{j}\delta_{jk},1\leq j,k\leq s\right\}  .
\end{align}%
\begin{equation}
\mathbf{\lambda}=\mathbf{\lambda}^{\dagger},\mathbf{\eta\in}\mathbb{R}^{s}%
\end{equation}%
\begin{align}
\mathcal{L}\left(  \mathbf{\varphi,s}\right)   &  =\mathcal{E}\left(
\mathbf{n}_{0}\mathbf{,\varphi}\right)  -\sum_{1\leq j,k\leq2s}\varepsilon
_{kj}\left(  \left\langle \left.  \varphi_{j}\right\vert \varphi
_{k}\right\rangle -s_{j}\delta_{jk}\right) \\
\frac{\partial\mathcal{L}\left(  \mathbf{\varphi,s}\right)  }{\partial s_{j}}
&  =\varepsilon_{jj}.
\end{align}%
\begin{equation}
\mathcal{L}\left(  \mathbf{n\left(  \mathbf{\eta}\right)  ,\varphi}\right)
=\mathcal{L}\left(  \mathbf{\varphi,s}\left(  \mathbf{\eta}\right)  \right)
\end{equation}


\section{Sensitivity}

Consider the optimization problem
\[
\min_{\mathbf{x\in}\Delta\left(  \mathbf{c}\right)  }f\left(  \mathbf{x}%
\right)  ,
\]
where
\[
\mathbf{\Delta}\left(  \mathbf{c}\right)  \mathcal{=}\left\{  \mathbf{x}%
\left\vert g\left(  \mathbf{x}\right)  =\mathbf{c}\right.  \right\}
\]
This can be solved using the Lagrangian
\[
\mathcal{L}\left(  \mathbf{x,\lambda}\right)  =f\left(  \mathbf{x}\right)
-\mathbf{\lambda}^{t}\left(  g\left(  \mathbf{x}\right)  -\mathbf{c}\right)
\]
$\left\{  \mathbf{x}_{m}\left(  \mathbf{\lambda}\left(  \mathbf{c}\right)
\right)  \right\}  $ are constrained stationary points if
\[
\nabla_{\mathbf{x}}\mathcal{L}\left(  \mathbf{x}_{m}\left(  \mathbf{\lambda
}\left(  \mathbf{c}\right)  \right)  \mathbf{,\lambda}\left(  \mathbf{c}%
\right)  \right)  =\nabla_{\mathbf{x}}f\left(  \mathbf{x}_{m}\left(
\mathbf{\lambda}\left(  \mathbf{c}\right)  \right)  \right)  -\mathbf{\lambda
}\left(  \mathbf{c}\right)  ^{t}\nabla_{\mathbf{x}}g\left(  \mathbf{x}%
_{m}\left(  \mathbf{\lambda}\left(  \mathbf{c}\right)  \right)  \right)
=\mathbf{0}%
\]
and
\[
g\left(  \mathbf{x}_{m}\left(  \mathbf{\lambda}\left(  \mathbf{c}\right)
\right)  \right)  -\mathbf{c=0}%
\]
for
\[
\mathbf{\lambda}\left(  \mathbf{c}\right)  \mathbf{\in\Lambda}\left(
\mathbf{c}\right)
\]
where $\mathbf{\Lambda}\left(  \mathbf{c}\right)  $ is the set of Lagrange
parameters for which eq. $\left(  \ref{sc1}\right)  $ has solutions. $\lceil
$The procedure is: find a $\lambda\left(  \mathbf{c}\right)  $ then a set of
solutions $\Gamma\left(  \lambda\left(  \mathbf{c}\right)  \right)  =\left\{
\mathbf{x}_{m}\left(  \lambda\left(  \mathbf{c}\right)  \right)  \right\}  $
for that $\lambda\left(  \mathbf{c}\right)  \rfloor$. The sets
$\mathbf{\Lambda}\left(  \mathbf{c}\right)  $ and $\Gamma\left(
\lambda\left(  \mathbf{c}\right)  \right)  $ may be empty.

The stationary points are minima if
\begin{equation}
\nabla_{\mathbf{x}}^{2}f\left(  \mathbf{x}_{m}\right)  \geq0\text{ in the
feasible subspace}%
\end{equation}
they are non degenerate minima if $\nabla_{\mathbf{x}}^{2}f\left(
\mathbf{x}_{m}\right)  $ is positive definite, but degenerate minima if
$\nabla_{\mathbf{x}}^{2}f\left(  \mathbf{x}_{m}\right)  $ is only positive.

Given an optimization problem
\[
f\left(  \mathbf{x}\left(  \mathbf{p}\right)  \right)  =\min_{\mathbf{x\in
\Delta}\left(  \mathbf{p}\right)  }f\left(  \mathbf{x}\right)
\]
where
\[
\Delta\left(  \mathbf{p}\right)  =\left\{  \left.  \mathbf{x}\right\vert
g\left(  \mathbf{x}\right)  =\mathbf{p}\right\}
\]
one can define a Lagrangian
\[
\mathcal{L}\left(  \mathbf{x,\mu,p}\right)  =f\left(  \mathbf{x}\right)
-\mathbf{\mu\bullet}\left(  g\left(  \mathbf{x}\right)  -\mathbf{p}\right)
\]
Stationary points satisfy the necessary first order conditions
\begin{align*}
\nabla_{\mathbf{x}}\mathcal{L}\left(  \mathbf{x}\left(  \mathbf{p}\right)
\mathbf{,\mu\left(  \mathbf{p}\right)  ,p}\right)   &  =0\\
\nabla_{\mathbf{\mu}}\mathcal{L}\left(  \mathbf{x\left(  \mathbf{p}\right)
,\mu\left(  \mathbf{p}\right)  ,p}\right)   &  =0
\end{align*}
We can define
\[
f\left(  \mathbf{x}\left(  \mathbf{p}\right)  \right)  =\min_{\mathbf{x\in
\Delta}\left(  \mathbf{p}\right)  }f\left(  \mathbf{x}\right)
\]%
\begin{equation}
\Omega=\left\{  \left.  \left(  \mathbf{x\left(  \mathbf{p}\right)
,\mu\left(  \mathbf{p}\right)  }\right)  \right\vert \right\}
\end{equation}


Under certain regularity conditions

$\mathbf{x\left(  \mathbf{p}\right)  }$ and $\mathbf{\mu\left(  \mathbf{p}%
\right)  }$ are functions of $\mathbf{p}$

and one can consider the restriction

of the Lagrangian and the object function 

to the space of optimal values
\[
\mathcal{L}\left(  \mathbf{x\left(  \mathbf{p}\right)  ,\mu\left(
\mathbf{p}\right)  }\right)  \equiv\mathcal{L}\left(  \mathbf{x}\left(
\mathbf{p}\right)  \mathbf{,\mu\left(  \mathbf{p}\right)  ,p}\right)
\]
and $f\left(  \mathbf{x}\left(  \mathbf{p}\right)  \right)  $. One can show
that
\[
\mathcal{L}\left(  \mathbf{x\left(  \mathbf{p}\right)  ,\mu\left(
\mathbf{p}\right)  }\right)  =f\left(  \mathbf{x}\left(  \mathbf{p}\right)
\right)
\]
The derivative of the Lagrangian wrt the constraint value $\mathbf{p}$ is
\[
\nabla_{\mathbf{p}}\mathcal{L}\left(  \mathbf{x,\mu,p}\right)  =\mathbf{\mu}%
\]
evaluating this derivative on the set $\mathbf{x}\left(  \mathbf{p}\right)  $
gives the sensitivity of the optimal solutions to changes in the constraint
\begin{align*}
\nabla_{\mathbf{p}}\mathcal{L}\left(  \mathbf{x\left(  \mathbf{p}\right)
,\mu\left(  \mathbf{p}\right)  }\right)   &  =\nabla_{\mathbf{p}}f\left(
\mathbf{x}\left(  \mathbf{p}\right)  \right)  \\
&  =\mathbf{\mu\left(  \mathbf{p}\right)  }%
\end{align*}


\subsection{Dual Problem}

\bigskip One can consider the Lagrangian $\mathcal{L}\left(  \mathbf{x,\lambda
}\right)  $ to be function of $\mathbf{x}$ and $\mathbf{\lambda}$ then the
stationarity conditions are
\begin{subequations}
\begin{align}
\nabla_{\mathbf{x}}\mathcal{L}\left(  \mathbf{x,\lambda}\right)   &
=\nabla_{\mathbf{x}}f\left(  \mathbf{x}\right)  -\mathbf{\lambda}^{t}%
\nabla_{\mathbf{x}}g\left(  \mathbf{x}\right)  =\mathbf{0}\label{d1}\\
\nabla_{\mathbf{\lambda}}\mathcal{L}\left(  \mathbf{x,\lambda}\right)   &
=-\left(  g\left(  \mathbf{x}\right)  -\mathbf{c}\right)  =\mathbf{0}
\label{d2}%
\end{align}
one should note that the condition eq. $\left(  \ref{d2}\right)  $ is the
constraint. We denote solutions to these equations as $\left\{  \left(
\mathbf{x}_{s}\left(  \mathbf{c}\right)  \mathbf{,\lambda}_{s}\left(
\mathbf{c}\right)  \right)  \right\}  $. The Hessian of the Lagrangian is
\end{subequations}
\begin{equation}
\left(
\begin{array}
[c]{cc}%
\nabla_{\mathbf{x}}^{2}\mathcal{L}\left(  \mathbf{x,\lambda}\right)  &
\nabla_{\mathbf{x\lambda}}^{2}\mathcal{L}\left(  \mathbf{x,\lambda}\right) \\
\nabla_{\mathbf{\lambda x}}^{2}\mathcal{L}\left(  \mathbf{x,\lambda}\right)  &
\nabla_{\mathbf{\lambda}}^{2}\mathcal{L}\left(  \mathbf{x,\lambda}\right)
\end{array}
\right)  ,
\end{equation}
where
\begin{align}
\nabla_{\mathbf{x\lambda}}^{2}\mathcal{L}\left(  \mathbf{x,\lambda}\right)
&  =-\nabla_{\mathbf{x}}g\left(  \mathbf{x}\right) \nonumber\\
\nabla_{\mathbf{\lambda x}}^{2}\mathcal{L}\left(  \mathbf{x,\lambda}\right)
&  =-\nabla_{\mathbf{x}}g\left(  \mathbf{x}\right)  ^{t}%
\end{align}
and
\begin{equation}
\nabla_{\mathbf{\lambda}}^{2}\mathcal{L}\left(  \mathbf{x,\lambda}\right)
=\mathbb{O}.
\end{equation}
If $\mathbf{x}_{m}\left(  \mathbf{c}\right)  $ is a local minimum then%
\begin{align}
&  \left(
\begin{array}
[c]{cc}%
\delta\mathbf{x} & \mathbf{0}%
\end{array}
\right)  \left(
\begin{array}
[c]{cc}%
\nabla_{\mathbf{x}}^{2}\mathcal{L}\left(  \mathbf{x}_{m}\left(  \mathbf{c}%
\right)  \mathbf{,\lambda}_{m}\left(  \mathbf{c}\right)  \right)  &
-\nabla_{\mathbf{x}}g\left(  \mathbf{x}_{m}\left(  \mathbf{c}\right)  \right)
\\
-\nabla_{\mathbf{x}}g\left(  \mathbf{x}_{m}\left(  \mathbf{c}\right)  \right)
^{t} & \mathbb{O}%
\end{array}
\right)  \left(
\begin{array}
[c]{c}%
\delta\mathbf{x}\\
\mathbf{0}%
\end{array}
\right)  \geq0\forall\delta\mathbf{x\backepsilon}\delta\mathbf{x}^{t}%
\nabla_{\mathbf{x}}g\left(  \mathbf{x}_{m}\left(  \mathbf{c}\right)  \right)
=0\nonumber\\
&  \qquad\qquad\qquad\qquad\qquad\qquad\qquad\qquad\qquad\Updownarrow
\nonumber\\
&  \qquad\qquad\qquad\qquad\qquad\qquad\delta\mathbf{x}^{t}\nabla_{\mathbf{x}%
}^{2}\mathcal{L}\left(  \mathbf{x}_{m}\left(  \mathbf{c}\right)
\mathbf{,\lambda}_{m}\left(  \mathbf{c}\right)  \right)  \delta\mathbf{x}%
\geq0\mathbf{\backepsilon}\delta\mathbf{x}^{t}\nabla_{\mathbf{x}}g\left(
\mathbf{x}_{m}\left(  \mathbf{c}\right)  \right)  =0
\end{align}
and $\mathbf{\lambda}_{m}\left(  \mathbf{c}\right)  $ is a local maximum
\begin{align}
&  \left(
\begin{array}
[c]{cc}%
\mathbf{0} & \delta\mathbf{\lambda}%
\end{array}
\right)  \left(
\begin{array}
[c]{cc}%
\nabla_{\mathbf{x}}^{2}\mathcal{L}\left(  \mathbf{x}_{m}\left(  \mathbf{c}%
\right)  \mathbf{,\lambda}_{m}\left(  \mathbf{c}\right)  \right)  &
\nabla_{\mathbf{x\lambda}}^{2}\mathcal{L}\left(  \mathbf{x}_{m}\left(
\mathbf{c}\right)  \mathbf{,\lambda}_{m}\left(  \mathbf{c}\right)  \right) \\
\nabla_{\mathbf{\lambda x}}^{2}\mathcal{L}\left(  \mathbf{x}_{m}\left(
\mathbf{c}\right)  \mathbf{,\lambda}_{m}\left(  \mathbf{c}\right)  \right)  &
\mathbb{O}%
\end{array}
\right)  \left(
\begin{array}
[c]{c}%
\mathbf{0}\\
\delta\mathbf{\lambda}%
\end{array}
\right)  \leq0\nonumber\\
&  \qquad\qquad\qquad\qquad\qquad\qquad\qquad\qquad\qquad\Updownarrow
\nonumber\\
&  \qquad\qquad\qquad\qquad\qquad\qquad\delta\mathbf{\lambda}^{t}%
\nabla_{\mathbf{x}}g\left(  \mathbf{x}_{m}\left(  \mathbf{c}\right)  \right)
\delta\mathbf{\lambda}\geq0,
\end{align}
i.e. $\left(  \mathbf{x}_{m}\left(  \mathbf{c}\right)  \mathbf{,\lambda}%
_{m}\left(  \mathbf{c}\right)  \right)  $ is a saddle point.

\bigskip%

\[
\frac{\delta\mathcal{E}\left(  \mathbb{D}^{2},\bm{\varphi}\right)  }%
{\delta\varphi_{j}^{\ast}}=\sum_{1\leq j,k\leq2s}\left\vert \varphi
_{k}\right\rangle \varepsilon_{kj}.
\]
We show in section \ref{Fock} that
\[
\frac{\delta\mathcal{E}\left(  \mathbb{D}^{2},\bm{\varphi}\right)  }%
{\delta\varphi_{j}^{\ast}}=F\left(  \mathbb{D}^{2},\bm{\varphi}\right)
\left\vert \varphi_{j}\right\rangle
\]
so eq.\ $\left(  \ref{statcon}\right)  $ becomes
\[
F\left(  \mathbb{D}^{2},\bm{\varphi}\right)  \left\vert \varphi_{j}%
\right\rangle =\sum_{1\leq j,k\leq2s}\left\vert \varphi_{k}\right\rangle
\varepsilon_{kj}.
\]%
\[
F\left(  \mathbb{D}^{2},\bm{\varphi}\right)  =F\left(  \mathbb{D}%
^{2},\bm{\varphi}\right)  ^{\dagger}%
\]
As (see section \ref{Fock})%
\[
\frac{\delta\mathcal{E}\left(  \bm{\varphi}\right)  }{\delta\varphi_{j}%
}=\left(  \frac{\delta\mathcal{E}\left(  \bm{\varphi}\right)  }{\delta
\varphi_{j}^{\ast}}\right)  ^{\ast}%
\]
\begin{subequations}
\label{Fock_Op}%
\begin{align*}
\qquad\qquad &  \qquad\qquad\qquad\qquad\qquad\left.  F\left(  \mathbb{D}%
^{2},\bm{\varphi}\right)  \right. \\
&  \qquad\qquad\qquad\qquad\qquad\qquad\left.  \shortparallel\right. \\
&  \sum_{1\leq m,n\leq2s}\left\{  \left(  \mathbb{H}_{1}\left(
\bm{\varphi}\right)  \mathbb{D}^{1}\right)  _{mn}+\frac{1}{2}\left(  L_{2}%
^{1}\left(  \mathbb{V}_{2}\left(  \bm{\varphi}\right)  \mathbb{D}^{2}\right)
\right)  _{mn}\right\}  \left\vert \varphi_{m}\right\rangle \left\langle
\varphi_{n}\right\vert
\end{align*}%
\end{subequations}
\[
F\left(  \mathbb{D}^{2},\bm{\varphi}\right)  ^{\dagger}=\sum_{1\leq m,n\leq
2s}\left\{  \left(  \mathbb{D}^{1}\mathbb{H}_{1}\left(  \bm{\varphi }\right)
\right)  _{mn}+\left(  L_{2}^{1}\left(  \mathbb{D}^{2}\mathbb{H}_{2}\left(
\bm{\varphi}\right)  \right)  \right)  _{mn}\right\}  \left\vert \varphi
_{m}\right\rangle \left\langle \varphi_{n}\right\vert .
\]


$I\left(  \psi_{p}\right)  $ depend on CGSO's and

the canonical coefficients%

\[
\left.  {}\right.
\]


$h_{1}\left(  \psi_{p}\right)  $ only depends on NGSO's and

occupation numbers

i.e. it is constant on $\left[  D^{1}\right]  _{N}$

\bigskip%

\begin{subequations}
\label{FE}%
\begin{align*}
\mathcal{E}\left(  \mathbb{D}^{2},\bm{\varphi}\right)   &  =\operatorname{Tr}%
\left\{  F\left(  \mathbb{D}^{2},\bm{\varphi}\right)  \right\}  -\mathcal{E}%
_{2}\left(  \mathbb{D}^{2},\bm{\varphi}\right) \\
\mathcal{E}\left(  \mathbb{D}^{2},\bm{\varphi}\right)   &  =\frac{1}%
{2}\left\{  \overset{}{\operatorname{Tr}\left\{  F\left(  \mathbb{D}%
^{2},\bm{\varphi}\right)  \right\}  +\mathcal{E}_{1}\left(  \mathbb{D}%
^{2},\bm{\varphi}\right)  }\right\}
\end{align*}%
\end{subequations}
\[
\mathcal{E}\left(  \mathbb{D}^{2},\bm{\varphi}\right)  =\frac{1}{2}\left\{
\overset{}{\operatorname{Tr}\left\{  F\left(  \mathbb{D}^{2}%
,\bm{\varphi}\right)  \right\}  +h_{1}\left(  \mathbb{D}^{2}%
,\bm{\varphi}\right)  }\right\}
\]


\bigskip%

\[
\varepsilon_{jk}=\varepsilon_{kj}^{\ast}%
\]


Thus we can express the optimization problem in two ways GSGS
\begin{align*}
\mathcal{E}\left(  \mathbf{\eta}_{GS},\mathbf{\psi}\left(  \mathbf{\eta}%
_{GS}\right)  \right)   &  =\min_{\mathbf{\eta\in}\mathbb{R}^{s}}\left\{
\min_{\mathbf{\psi\in\Delta}}\mathcal{E}\left(  \mathbf{\eta,\psi}\right)
\right\} \\
&  =\min_{\mathbf{\eta\in}\mathbb{R}^{s}}\mathcal{E}\left(  \mathbf{\eta,\psi
}\left(  \mathbf{\eta}\right)  \right)
\end{align*}%
\[
\mathcal{E}\left(  \mathbf{\eta,\psi}\left(  \mathbf{\eta}\right)  \right)
=\min_{\mathbf{\psi\in\Delta}}\mathcal{E}\left(  \mathbf{\eta,\psi}\right)
\]
where
\[
\mathbf{\Delta}=\left\{  \left\vert \psi_{j}\right\rangle \left\vert
\overset{}{\left\langle \left.  \psi_{j}\right\vert \psi_{k}\right\rangle
=\delta_{jk};\left\vert \psi_{j}\right\rangle \in\mathcal{H}^{1}}\right.
\right\}
\]%
\[
\mathbf{\eta\in}\mathbb{R}^{s}%
\]
and
\[
\mathcal{E}\left(  \mathbf{\Lambda}_{GS}\right)  =\min_{\mathbf{\Lambda\in
}\mathbb{R}^{s}\times\mathbb{R}^{u}\times\mathbb{R}^{s}}\mathcal{E}\left(
\mathbf{\Lambda}\right)
\]
where
\[
\mathbf{\Lambda}=\left(  \mathbf{\eta,\alpha}\left(  \mathbf{x,y}\right)
,\mathbf{\theta}\right)  ;\quad\mathbf{\eta\in}\mathbb{R}^{s},\mathbf{x,y\in
}\mathbb{R}^{u},\mathbf{\theta\in}\mathbb{R}^{s}%
\]%
\[
\mathbf{\alpha}=\mathbf{x}+i\mathbf{y}%
\]%
\[
\mathbf{\eta\in}\mathbb{R}^{s},\mathbf{x,y\in}\mathbb{R}^{u},\mathbf{\theta
\in}\mathbb{R}^{s}%
\]
One can define
\[
\Omega=\left\{  \left.  \mathbf{\Upsilon}_{GS}\left(  \mathbf{\eta}\right)
=\min_{\mathbf{\Upsilon\in\Delta}\left(  \mathbf{\eta}\right)  }%
\mathcal{E}\left(  \mathbf{\Upsilon}\right)  \right\vert \mathbf{\eta\in
}\mathbb{R}^{s}\right\}
\]
so that the ground energy is given by
\[
\mathcal{E}\left(  \mathbf{\Upsilon}_{GS}\left(  \mathbf{\eta}_{\ast}\right)
\right)  =\min_{\mathbf{\eta\in}\mathbb{R}^{s}}\mathcal{E}\left(
\mathbf{\Upsilon}_{GS}\left(  \mathbf{\eta}\right)  \right)
\]
It is possible that for a given $\mathbf{\eta}$ the minimum is degenerate so
one obtains solutions $\left\{  \mathbf{\Upsilon}_{GS1}\left(  \mathbf{\eta
}\right)  ,\cdots,\mathbf{\Upsilon}_{GSm}\left(  \mathbf{\eta}\right)
\right\}  ,$ which also must be included in the set $\Omega$. One can define a
Lagrangian
\[
\mathcal{L}\left(  \mathbf{\Upsilon}\right)  =\mathcal{E}\left(
\mathbf{\Upsilon}\right)  -\sum_{1\leq j,k\leq2s}\varepsilon_{jk}\left(
\left\langle \left.  \Upsilon_{j}\right\vert \Upsilon_{k}\right\rangle
-c_{0j\mathit{s}\mathbf{\eta}}\mathit{s}\left(  \mathbf{\eta}\right)
^{-1}e^{\eta_{j}}\delta_{jk}\right)
\]
associated with the constrained variation (orthogonal with a fixed
normalization) of the energy functional $\mathcal{E}\left(  \mathbf{\Upsilon
}\right)  ,$ where
\[
\mathcal{L}:\overset{s}{\overbrace{\mathcal{H}^{1}\times\cdots\times
\mathcal{H}^{1}}}\times\,\mathbb{R\times R}\rightarrow\mathbb{R}%
\]
The stationarity conditions are
\[
\frac{\delta\mathcal{L}\left(  \mathbf{\Upsilon}\left(  \mathbf{c}\right)
\right)  }{\delta\Upsilon_{j}^{\ast}}=\frac{\delta\mathcal{E}\left(
\mathbf{\Upsilon}\left(  \mathbf{c}\right)  \right)  }{\delta\Upsilon
_{j}^{\ast}}-\sum_{1\leq k\leq2s}\varsigma_{jk}\left(  \mathbf{c}\right)
\left\vert \Upsilon_{k}\left(  \mathbf{c}\right)  \right\rangle =0
\]


\section{}
\end{document}